\chapter{数据通路与控制器}

\begin{keyidea}
一条指令必须落实为源状态、组合路径、控制线和提交边界。本章把 ALU、寄存器、选择器、控制字和状态机接成可走读的数据通路。
\end{keyidea}

\section{一、数据通路让值从旧状态走到新状态}

数据通路回答一个问题：值从哪里来，经过哪里，最后保存到哪里。孤立的寄存器和 ALU 还不能完成连续动作。必须把寄存器阵列、选择器、ALU、内部连线和写回路径接成一个可控制的闭环。若 ALU 是“能算什么”，数据通路就是“值能走哪条路”。

一个寄存器只能保存一个值。若硬件需要同时保存多个中间值，就需要寄存器阵列。寄存器阵列可以理解为一组编号寄存器，加上读写选择电路。

\begin{longtable}{L{0.22\linewidth}L{0.42\linewidth}L{0.24\linewidth}}
\toprule
端口 & 作用 & 类型 \\
\midrule
读地址 A/B & 选择两个源寄存器 & 输入 \\\tablerowrule
读数据 A/B & 输出两个源值 & 输出 \\\tablerowrule
写地址 & 选择目标寄存器 & 输入 \\\tablerowrule
写数据 & 提供要保存的新值 & 输入 \\\tablerowrule
写使能 & 决定是否在时钟沿写入 & 输入 \\
\bottomrule
\end{longtable}

读地址进入译码和选择网络，决定哪两个寄存器的值出现在读数据端口。写地址进入译码器，写使能有效时，只允许被选中的那个寄存器在时钟沿接收写数据。这样，一组寄存器就能被编号访问，而不是每个寄存器单独拉一套控制线。真实处理器中的寄存器堆常常有多个读端口和写端口，原因也很直接：ALU 常常需要同时读两个源操作数，并在周期末写回一个目标。

\begin{pdfdiagram}
  \node[hw control, minimum width=2.4cm] (dec) at (0,0) {写地址译码\\+ \code{reg\_write}};
  \node[hw state, minimum width=1.4cm] (r0) at (3.4,1.2) {R0};
  \node[hw state, minimum width=1.4cm] (r1) at (3.4,0.4) {R1};
  \node[hw state, minimum width=1.4cm] (r2) at (3.4,-0.4) {R2};
  \node[hw state, minimum width=1.4cm] (r3) at (3.4,-1.2) {R3};
  \node[hw compute, minimum width=2.0cm] (ma) at (7.1,0.8) {读 mux A};
  \node[hw compute, minimum width=2.0cm] (mb) at (7.1,-0.8) {读 mux B};
  \draw[hw bus] (dec.east) -- (1.5,0) -- (1.5,1.2) -- (r0.west);
  \draw[hw bus] (1.5,0.4) -- (r1.west);
  \draw[hw bus] (1.5,-0.4) -- (r2.west);
  \draw[hw bus] (1.5,-1.2) -- (r3.west);
  \node[hw label, anchor=east] at (-2.4,-1.9) {写数据};
  \draw[hw bus] (-2.4,-1.6) -- (2.2,-1.6) -- (2.2,1.2) -- (r0.south);
  \draw[hw bus] (2.2,0.4) -- (r1.south);
  \draw[hw bus] (2.2,-0.4) -- (r2.south);
  \draw[hw bus] (2.2,-1.2) -- (r3.south);
  \draw[hw bus] (r0.east) -- (5.2,1.2);
  \draw[hw bus] (r1.east) -- (5.2,0.4);
  \draw[hw bus] (r2.east) -- (5.2,-0.4);
  \draw[hw bus] (r3.east) -- (5.2,-1.2);
  \draw[hw bus] (5.2,1.2) -- (5.2,-1.2);
  \draw[fill=BookInk] (5.2,1.2) circle (0.05);
  \draw[fill=BookInk] (5.2,0.4) circle (0.05);
  \draw[fill=BookInk] (5.2,-0.4) circle (0.05);
  \draw[fill=BookInk] (5.2,-1.2) circle (0.05);
  \draw[hw bus] (5.2,0.8) -- (ma.west);
  \draw[hw bus] (5.2,-0.8) -- (mb.west);
  \node[hw label] at (5.2,1.75) {四个候选值都进入两个读 mux};
  \node[hw label] at (6.5,2.1) {读地址 A};
  \draw[hw bus] (6.9,1.85) -- (ma.north);
  \node[hw label] at (6.5,-2.1) {读地址 B};
  \draw[hw bus] (6.9,-1.85) -- (mb.south);
  \node[hw label] at (10.1,0.8) {读数据 A};
  \node[hw label] at (10.1,-0.8) {读数据 B};
  \draw[hw bus] (ma.east) -- (9.4,0.8);
  \draw[hw bus] (mb.east) -- (9.4,-0.8);
\end{pdfdiagram}
\diagramnote{4 寄存器堆的内部结构：写地址经译码（并和 \code{reg\_write} 相与）只打开一个目标的写入；写数据广播到所有寄存器的 D 端，但只有被选者提交；两个读 mux 按读地址各选一路输出。读是组合选择，写是时钟沿提交——这条不对称是寄存器堆和“一组变量”的本质区别。}

ALU 的两个输入不一定总来自同一处。一个输入可能来自寄存器阵列，也可能来自 PC、计数器或固定 0；另一个输入可能来自寄存器、常量或扩展后的字段。因此 ALU 输入前通常有 mux。

\begin{lstlisting}
alu_a = mux(a_sel, reg_a, pc, zero, latch_a)
alu_b = mux(b_sel, reg_b, immediate, one, offset)
\end{lstlisting}

以寄存器 R1 和 R2 相加、结果写入 R3 为例。数据通路需要让读地址 A 选择 R1，读地址 B 选择 R2；\code{a_sel} 选择 \code{reg_a}，\code{b_sel} 选择 \code{reg_b}；\code{alu_op} 选择 ADD。此时 ALU 输出就是 R1+R2。

写回数据也不一定只来自 ALU。它可能来自 ALU，也可能来自存储器读口、外部输入、PC+4（字节编址惯例下，下一条 32-bit 指令的地址）或专用状态寄存器。写回前也需要 mux。

\begin{lstlisting}
write_data = mux(wb_sel, alu_result,
                 memory_data, input_data, pc_plus_4)
\end{lstlisting}

在 \code{R3=R1+R2} 的动作中，\code{wb_sel} 选择 \code{alu_result}，写地址选择 R3，\code{reg_write} 为 1。等时钟沿到来，R3 才真正保存结果。若 \code{reg_write} 为 0，即使 ALU 算出了正确结果，寄存器阵列也不会改变。这个边界来自“状态、时钟与时序逻辑”一章的同步设计原则：组合逻辑可以在周期内传播和抖动，状态只在受控时钟沿提交。

当多个部件都可能把值送到同一个目的地时，不能让它们同时强行驱动同一条线。要么用 mux 明确选择一路，要么用总线协议规定同一时刻只有一个驱动者。ALU 输出、存储器输出和外部输入如果直接接在一起，一旦两个源同时输出不同电平，就会形成冲突。mux 让这件事变成明确选择：当前周期只允许一路成为 \code{write_data}，适合芯片内部数据通路和 FPGA 逻辑，代价是输入多时选择器延迟和面积增加。三态总线也能共享连线，适合板级总线和部分分立芯片实验，但必须保证同一时刻只有一个输出使能有效，多个驱动同时打开就是冲突。点到点连线适合高频、固定方向的数据路径，代价是连线数量增加、灵活性降低。在初学实验和 FPGA 设计中，优先使用 mux 更容易检查和调试。

\topic{三态总线型数据通路：一条线上轮流说话}

“数字硬件基础”部分讲三态门时说过：三态输出除了 0 和 1，还有高阻态 Z——输出端像被“拔掉”一样，既不驱高也不驱低。三态总线正是利用这一点组织数据通路：把多个寄存器的输出、ALU 的输出、存储器读出都接到同一组公共线上，每个源前面放一个三态缓冲器，各自由输出使能控制。任何时刻最多一个使能有效，总线上就只有这一家在“说话”，其余源都处在高阻态，对总线电平没有影响。写入侧则反过来：总线上的值广播到所有寄存器的输入端，写使能决定谁在时钟沿接收。

把这段话画成数据通路，就是下面的结构：所有源挂到同一条总线上，读和写各由一组使能控制。

\begin{pdfdiagram}
  \node[hw state, minimum width=1.7cm, align=center] (ra) at (0.9,2.6) {寄存器 A};
  \node[hw state, minimum width=1.7cm, align=center] (rb) at (3.3,2.6) {寄存器 B};
  \node[hw compute, minimum width=1.7cm, align=center] (alu) at (5.7,2.6) {ALU 输出};
  \node[hw memory, minimum width=1.9cm, align=center] (mem) at (8.1,2.6) {存储器读出};
  \node[hw compute, minimum width=1.1cm, align=center] (ta) at (0.9,1.35) {三态\\门};
  \node[hw compute, minimum width=1.1cm, align=center] (tb) at (3.3,1.35) {三态\\门};
  \node[hw compute, minimum width=1.1cm, align=center] (talu) at (5.7,1.35) {三态\\门};
  \node[hw compute, minimum width=1.1cm, align=center] (tm) at (8.1,1.35) {三态\\门};
  \draw[BookMuted, line width=1.6pt] (0.1,0) -- (8.9,0);
  \node[hw label, anchor=east] at (-0.05,0) {公共数据总线};
  \draw[hw bus] (ra.south) -- (ta.north);
  \draw[hw bus] (ta.south) -- (0.9,0);
  \draw[hw bus] (rb.south) -- (tb.north);
  \draw[hw bus] (tb.south) -- (3.3,0);
  \draw[hw bus] (alu.south) -- (talu.north);
  \draw[hw bus] (talu.south) -- (5.7,0);
  \draw[hw bus] (mem.south) -- (tm.north);
  \draw[hw bus] (tm.south) -- (8.1,0);
  \node[hw control, minimum width=4.6cm, align=center] (ctl) at (4.5,-1.6) {互斥使能译码\\任一时刻只有一个 OE 有效};
  \draw[hw return, dashed] (ctl.north) -- (4.5,-0.7);
  \draw[draw=BookTeal, line width=0.58pt, dashed] (1.8,-0.7) -- (9.0,-0.7);
  \draw[hw return, dashed] (1.8,-0.7) -- (1.8,1.35) -- (ta.east);
  \draw[hw return, dashed] (4.2,-0.7) -- (4.2,1.35) -- (tb.east);
  \draw[hw return, dashed] (6.6,-0.7) -- (6.6,1.35) -- (talu.east);
  \draw[hw return, dashed] (9.0,-0.7) -- (9.0,1.35) -- (tm.east);
  \node[hw state, text width=5.2cm, align=center] (wr) at (4.5,-3.4) {各寄存器写入口：总线值广播到全部输入端，写使能决定谁在时钟沿接收};
  \draw[hw bus] (8.6,0) -- (8.6,-3.4) -- (wr.east);
  \node[hw label, anchor=east] at (8.45,-2.6) {广播写回};
\end{pdfdiagram}
\diagramnote{三态总线型数据通路：寄存器 A/B、ALU、存储器的输出各经一个三态门挂到同一条公共总线上，使能译码保证任一时刻只有一个 OE 有效，其余源处于高阻态；总线值广播到全部寄存器写入口，写使能决定谁在时钟沿接收。“谁在驱动”由控制逻辑负责，两个使能同时有效就是总线冲突。}

与 mux 型数据通路对照，差别在“选择发生在哪里”。mux 方案里每条数据通路独立布线，由接收端的选择器挑一路——选择是集中的、点对点的，几路输入同时驱动各自的线互不相干。总线方案里，选择分散在每个驱动者的输出使能上，依赖一条全局约定：同一时刻只有一个驱动者。这条约定一旦被破坏，两个源同时驱动总线，一个驱高一个驱低，线上就是不可预期的电平，还可能流过过大的电流。

\begin{longtable}{L{0.2\linewidth}L{0.38\linewidth}L{0.32\linewidth}}
\toprule
对照点 & 三态总线型 & mux 型 \\
\midrule
连线数量 & 每个位宽只有一组公共线，源再多也不加线 & 每个源到每个接收点独立走线，随源数增长 \\\tablerowrule
面积代价 & 省布线、省选择器，适合布线紧张的场合 & 多路 mux 本身占面积，输入越多级数越深 \\\tablerowrule
出错模式 & 使能冲突（多驱）或全高阻（无驱），影响是全局性的 & 选择信号错即选错路，故障局限在单点 \\\tablerowrule
时序分析 & 总线电平依赖各使能的时序，要分析所有潜在驱动者 & 路径固定，逐条路径算延迟即可 \\\tablerowrule
调试入口 & 先查“此刻谁在驱动” & 直接观察选择信号和对应路径 \\
\bottomrule
\end{longtable}

早期小型机和微处理器爱用总线型，原因很实际：当时芯片内部布线层数少、布线资源紧张，分立搭建时电路板上的走线和插接件更是寸土寸金；一条公共总线让寄存器、ALU、存储器、I/O 都挂上来，扩展一个新部件只需再挂一个三态驱动器，不必为每个新源重铺一组线。PDP-11 的 Unibus、各类 8-bit 微处理器的地址/数据总线，都是这个思路。代价是把“谁在驱动”的正确性推给了控制设计：使能信号必须由译码逻辑保证互斥，切换驱动者时最好留一拍全高阻的间隔，避免旧驱动尚未关闭、新驱动已经打开的瞬间冲突。正因为这条约定难以形式化检查，现代芯片内部数据通路大多回到 mux 型；三态总线更多留在板级互连和存储器接口这类“真正共享一条物理线”的地方。

一个最小数据通路闭环可以写成：

\begin{lstlisting}
register file -> mux -> ALU -> mux -> register file
\end{lstlisting}

\begin{pdfdiagram}
  \node[hw state, minimum width=2.15cm] (rf) at (0,0) {寄存器阵列};
  \node[hw compute, minimum width=1.75cm] (muxa) at (2.6,0.55) {A MUX};
  \node[hw compute, minimum width=1.75cm] (muxb) at (2.6,-0.55) {B MUX};
  \node[hw compute, minimum width=1.85cm] (alu) at (5.25,0) {ALU};
  \node[hw compute, minimum width=1.95cm] (wb) at (7.85,0) {写回 MUX};
  \node[hw state, minimum width=1.8cm] (flags) at (5.75,-1.45) {FLAGS};
  \node[hw control, minimum width=2.15cm] (ctrl) at (2.6,-2.12) {控制线};
  \draw[hw bus] (rf.east) -- ++(0.55,0.55) -- (muxa.west);
  \draw[hw bus] (rf.east) -- ++(0.55,-0.55) -- (muxb.west);
  \draw[hw bus] (muxa.east) -- (alu.west);
  \draw[hw bus] (muxb.east) -- (alu.west);
  \draw[hw bus] (alu) -- (wb);
  \node[hw label] at (6.55,0.45) {result};
  \draw[hw bus] (alu.south) |- (flags.west);
  \draw[hw return] (wb.east) -- ++(0.55,0) |- (rf.north);
  \node[hw label] at (7.75,0.95) {write data};
  \draw[hw bus] (ctrl.north) -- (muxa.south);
  \draw[hw bus] (ctrl.north) -- (muxb.south);
  \draw[hw bus] (ctrl.east) -- ++(1.0,0) |- (alu.south);
  \draw[hw bus] (ctrl.east) -- ++(4.25,0) -| (wb.south);
  \draw[hw bus] (ctrl.west) -| (rf.south);
  \node[hw label, text width=8.9cm] at (4.1,-2.95) {数据通路图要区分“值走哪条路”和“控制线选择哪条路”。写回箭头绕到上方，避免穿过 ALU 和标签。};
\end{pdfdiagram}
\diagramnote{最小数据通路闭环。源值从寄存器阵列读出，经过输入选择和 ALU，最后由写回选择器在时钟边界写回。}

它能在一个周期内读取旧值、选择输入、完成组合计算，并在时钟沿保存结果。这个闭环已经能表达很多动作，但它还不会自己决定“当前周期做什么”。哪些地址送入寄存器阵列，ALU 做加法还是减法，写回选择哪一路，写使能是否打开，都需要控制器统一产生。

如果把这个闭环搭成分立电路，器件可以按功能分组选择。74HC574 或 74HC173 可以保存寄存器值；74HC138 可以把写地址译码成某个寄存器的写使能；74HC157 可以选择 ALU 输入或写回来源；74HC283/74HC86 组合成加减核心；LED 或逻辑分析仪用于观察寄存器输出、ALU 输出和写使能。这个实验的关键不是一次做完整指令集，而是让 \code{R1}、\code{R2}、\code{R3} 的值在时钟沿前后清晰可见。

这个 4 寄存器堆有一个重要限制：若用普通 74HC574，它虽有三态输出使能 \code{/OE}，但没有专门的写使能逻辑，就要额外设计时钟使能或 D 端保持路径；若用 74HC173 这类带使能和输出控制的寄存器，实验更容易分阶段观察。无论选哪种器件，原则相同：写地址只决定“谁有资格写”，全局 \code{reg_write} 决定“本周期是否写”，时钟沿决定“何时真正写”。读 mux 是组合选择，写入是时序提交，两者不能混在一起。

把各连接点的器件和检查要点再落实一层：R0 到 R3 的存储用 74HC173（4-bit 寄存器）或 74HC574（8-bit 成组保存），复位后初值确定、时钟沿前后值可观察；写地址译码用 74HC138 把两位 \code{write_addr} 展开成四路写使能候选，同一周期只允许一个目标写入；全局 \code{reg_write} 与译码输出相与后送入各寄存器，\code{reg_write=0} 时所有寄存器保持；读端口 A/B 各用一组 4-bit mux 按 \code{read_addr_a}、\code{read_addr_b} 独立选择，A 端口变化不应改写寄存器内容，A/B 可同时读同一寄存器或不同寄存器；写回数据经写回 mux 进入所有寄存器 D 端，未选寄存器即使 D 端变化也不能写入。

真实芯片把这些部件压缩进同一颗硅片。8051 把 CPU、片内 RAM、特殊功能寄存器、定时器、串口和 I/O 口组织在一个单片机内部；对程序员来说，某些外设像“寄存器地址”一样被读写；对硬件来说，这是数据通路、地址译码、外设寄存器和控制器共同工作的结果。8086 把总线接口单元和执行单元分工，执行单元内部仍要围绕寄存器、ALU、标志和控制路径组织动作。80386 继续扩大寄存器和地址计算规模，并引入更复杂的保护与寻址硬件。规模不同，问题相同：值必须从旧状态沿受控路径走到新状态。

\begin{chipcase}
从 8051 看，单片机不能简化为“小型软件平台”，而是 CPU 数据通路、片内存储器、I/O 口、定时器和中断控制的硬件组合。“数字硬件基础”部分的输入调理和逻辑门，“状态、时钟与时序逻辑”一章的状态边界，本章的寄存器与控制线，都在单片机内部或引脚边界重新出现。
\end{chipcase}

数据通路的时序分析应至少包含四项：源状态何时稳定，组合路径最坏延迟是多少，目标触发器建立/保持时间是否满足，哪些写使能在时钟沿有效。若一个单周期 \code{R3=R1+R2} 路径太长，可以把动作拆成多周期：第一周期读 R1/R2 并锁存到 A/B，第二周期 ALU 计算并锁存到 ALUOut，第三周期把 ALUOut 写回 R3。这样周期时间可以变短，但一次动作需要更多周期。速度、面积和控制复杂度再次形成取舍。

三种方案的取舍：单周期让读寄存器、ALU 计算、写回都在一个周期内完成，控制简单，但周期必须容纳最长路径；多周期把读、算、写分到多个状态和多个时钟沿，周期可短，代价是控制状态增加；流水线让多条动作在不同阶段重叠推进，吞吐提高，代价是冒险、旁路和冲刷更复杂。

\section{二、控制器输出一组一致的控制线}

数据通路提供很多可能路径，但同一周期不能所有路径都随意打开。控制器的任务是根据当前状态、指令字段和条件输入，产生一组控制信号，让数据通路执行一个明确动作。控制器不是写在纸上的命令列表，而是一块输出控制电线的硬件。它可以由组合逻辑、状态机、ROM、PLA（可编程逻辑阵列：把若干 AND 项和 OR 项做成规则阵列的组合逻辑，适合把译码表直接烧进硬件）、微码表或这些方法的混合实现。

控制信号通常是一组 0/1 或多 bit 选择值：

\begin{lstlisting}
alu_op
a_sel
b_sel
wb_sel
read_addr_a
read_addr_b
write_addr
reg_write
flags_write
mem_read
mem_write
next_state_sel
\end{lstlisting}

这些信号直接接到 ALU、mux、寄存器阵列、标志寄存器和存储器接口。控制信号给错，数据通路就会走错。例如 \code{reg_write} 本应为 0 却变成 1，某个寄存器会在时钟沿被意外覆盖；\code{wb_sel} 本应选择 ALU，却选了外部输入，写回值就会错；\code{flags_write} 本应关闭却打开，旧标志会被无关结果污染，后续条件分支会误判。

若动作规则简单，可以用组合逻辑直接从输入生成控制信号。例如输入 \code{mode} 为 0 时做加法，为 1 时做按位 AND：

\begin{lstlisting}
mode = 0 -> alu_op = ADD
mode = 1 -> alu_op = AND
\end{lstlisting}

硬连线控制速度快，结构也直接。缺点是规则复杂后，组合逻辑会变大，修改一个动作可能影响多条控制线。若一个动作无法在一个周期内完成，就需要状态机。状态机保存“当前走到哪一步”，每个状态输出一组控制信号，并决定下一状态。

控制器可以分成两个层次理解。第一层是\hwkey{译码}：把 opcode、功能位、条件位和当前状态翻译成控制线。第二层是\hwevidence{时序安排}：决定这些控制线在哪个周期有效，哪些状态元件在该周期写入。若指令简单、动作少，硬连线译码可以直接生成控制信号；若指令复杂、步骤多，控制器可能使用微操作表或微码 ROM，把一条指令拆成多行控制字逐步执行。

几种控制方式的取舍：硬连线控制适合动作规则固定、追求低延迟的路径，代价是逻辑复杂后修改困难、一个错误影响多条控制线；状态机控制适合多周期动作和明确阶段推进，代价是状态编码、非法状态和输出毛刺都要处理；微码控制适合指令复杂、步骤可表格化的场景，代价是微码存储、入口选择和异常中断点更复杂；混合控制让常见简单路径走硬连线、复杂路径走微操作，代价是必须统一提交边界和异常恢复规则。

一行控制字可以像下面这样理解：

\begin{lstlisting}
// 注释（推进）：当前状态与输入共同选择动作和下一状态，先确认每个分支的优先级。
// 注释（边界）：本拍只计算 next，状态在时钟边界更新；等待和错误都要有去向。
state EXEC_ADD:
  read_addr_a  = rs1
  read_addr_b  = rs2
  a_sel        = REG_A
  b_sel        = REG_B
  alu_op       = ADD
  wb_sel       = ALU_RESULT
  write_addr   = rd
  reg_write    = 1
  flags_write  = 1
  next_state   = FETCH
\end{lstlisting}

这并非软件语句，而是对硬件控制线的命名描述。控制字中的每一位最终都要连接到 mux、ALU、寄存器写使能、标志寄存器、存储器接口或下一状态逻辑。调试控制器时，不应只看“当前指令是什么”，还要看当前状态、控制字、条件输入和本周期会提交哪些状态。

控制器也可以用分立芯片做出教学模型。一个 74HC161 计数器可以充当微步骤计数器，74HC138 可以把微步骤译码成若干状态脉冲，小容量 EEPROM 或 ROM 可以存放控制字。地址线接 opcode、微步骤和标志位，数据线输出 \code{alu_op}、\code{a_sel}、\code{b_sel}、\code{reg_write} 等控制信号。这种做法的好处是非常直观：修改控制表就能改变机器动作；代价是控制存储器访问延迟、控制字宽度和异常入口都要进入时序设计。

\begin{pdfdiagram}
  \node[hw state, minimum width=2.4cm] (ir) at (0,1.35) {IR\\opcode[3:0]};
  \node[hw state, minimum width=2.4cm] (step) at (0,0.45) {微步骤计数器\\step[2:0]};
  \node[hw state, minimum width=2.4cm] (flg) at (0,-0.45) {标志寄存器\\Z/C/V/N};
  \node[hw state, minimum width=2.4cm] (rst) at (0,-1.35) {reset/irq};
  \node[hw memory, minimum width=3.3cm, minimum height=3.6cm] (rom) at (4.8,0) {控制 ROM\\地址 $\rightarrow$ 控制字};
  \node[hw compute, minimum width=3.5cm, minimum height=3.6cm, align=left] (out) at (10.4,0) {控制字字段\\\code{alu\_op / a\_sel / b\_sel}\\\code{wb\_sel / reg\_write}\\\code{flags\_write / pc\_write}\\\code{ir\_write / mem\_rw}\\\code{next\_sel}};
  \draw[hw bus] (ir.east) -- (rom.west |- ir.east);
  \draw[hw bus] (step.east) -- (rom.west |- step.east);
  \draw[hw bus] (flg.east) -- (rom.west |- flg.east);
  \draw[hw bus] (rst.east) -- (rom.west |- rst.east);
  \draw[hw bus] (rom.east) -- node[hw label, above]{控制字} (out.west);
  \draw[hw bus] (out.east) -- node[hw label, above]{到数据通路各控制点} ++(2.3,0);
\end{pdfdiagram}
\diagramnote{控制 ROM 的结构：指令 opcode、微步骤、标志位和复位/中断请求拼成 ROM 地址，地址内容就是本周期全组控制线。改机器动作等于改 ROM 内容表——这就是微码控制“指令是可编程表格”的含义。}

\begin{longtable}{L{0.2\linewidth}L{0.38\linewidth}L{0.32\linewidth}}
\toprule
控制 ROM 地址位 & 来源 & 作用 \\
\midrule
\code{opcode[3:0]} & 指令寄存器中的操作码 & 选择当前指令类型，如 ADD、SUB、LOAD、BRANCH \\\tablerowrule
\code{step[2:0]} & 微步骤计数器或状态寄存器 & 区分取指、译码、执行、写回等阶段 \\\tablerowrule
\code{Z/C/V/N} & 标志寄存器输出 & 条件分支或条件写回的输入 \\\tablerowrule
\code{reset/irq} & 复位或外部请求同步后的状态 & 强制进入复位入口或异常入口 \\
\bottomrule
\end{longtable}

控制 ROM 的数据输出就是控制字。为了便于搭实验机，可以先把控制字设计成固定宽度字段，而不是零散电线。

\begin{longtable}{L{0.16\linewidth}L{0.32\linewidth}L{0.42\linewidth}}
\toprule
控制字段 & 连接对象 & 含义 \\
\midrule
\code{alu_op[2:0]} & ALU 函数选择 & 指定 ADD、SUB、AND、OR、XOR、PASS 等动作 \\\tablerowrule
\code{a_sel[1:0]} & ALU A 输入 mux & 选择寄存器 A、PC、0 或临时寄存器 \\\tablerowrule
\code{b_sel[1:0]} & ALU B 输入 mux & 选择寄存器 B、立即数、1 或偏移 \\\tablerowrule
\code{wb_sel[1:0]} & 写回 mux & 选择 ALUOut、存储器数据、PC+1 或外部输入（PC+1 按指令编址惯例，即下一条指令）\\\tablerowrule
\code{reg_write} & 寄存器堆写使能 & 决定目标寄存器是否在时钟沿写入 \\\tablerowrule
\code{flags_write} & 标志寄存器写使能 & 决定 ALU 标志是否提交 \\\tablerowrule
\code{pc_write} & PC 写使能 & 决定下一 PC 是否提交 \\\tablerowrule
\code{ir_write} & 指令寄存器写使能 & 决定本周期是否锁存取到的指令 \\\tablerowrule
\code{mem_read/write} & 存储器接口 & 发起读周期或写周期 \\\tablerowrule
\code{next_sel[1:0]} & 下一状态逻辑 & 选择顺序微步骤、取指、分支目标或异常入口 \\
\bottomrule
\end{longtable}

这个表可以直接变成 EEPROM 内容表。例如 \code{opcode=ADD, step=EXEC} 的控制字会选择寄存器 A/B 作为 ALU 输入，\code{alu_op=ADD}，打开 \code{flags_write}，并把下一状态设为写回；\code{opcode=BRANCH, step=EXEC} 的控制字则可能关闭 \code{reg_write}，读取 Zero 标志，通过 \code{next_sel} 选择顺序 PC 或分支目标。判断控制 ROM 方案是否可行，关键不在“表看起来合理”，而在每一位输出都能追到一个实际器件引脚。

\topic{微程序控制器：把控制字存进一块存储器}

前面控制 ROM 的做法其实已经触到了微程序控制（microprogrammed control）的核心：与其用组合逻辑从 opcode 现场译出控制线，不如把“每条指令在每个周期该输出什么”事先写成一张表，存进一块专门的控制存储器。表里的每一行叫一条微指令，一台机器全部微指令合起来叫微程序。一条机器指令不再对应一段组合逻辑，而对应控制存储器里的一段微指令序列——取指之后跳到这段微程序，逐条执行，走完这段微程序，这条机器指令的所有周期也就走完了。

一条微指令通常包含三类字段：

\begin{longtable}{L{0.2\linewidth}L{0.42\linewidth}L{0.28\linewidth}}
\toprule
字段 & 作用 & 例子 \\
\midrule
控制线字段 & 本周期送往数据通路的全部控制位，每位接一条控制线 & \code{alu_op}、\code{a_sel}、\code{reg_write}、\code{mem_read} \\\tablerowrule
下一微地址字段 & 指出正常情况下下一条微指令的地址 & \code{next =} 微地址 17 \\\tablerowrule
条件选择字段 & 指定本周期是否查看条件、查看哪个条件，条件成立时改走哪个地址 & 看 Zero：成立去分支微地址，不成立用下一微地址 \\
\bottomrule
\end{longtable}

负责读这张表的小硬件叫微序列器（microsequencer）。它每个周期做同一件事：按当前微地址取出一条微指令，把控制线字段直接送往数据通路，同时按条件选择字段决定下一个微地址——顺序、跳到指定地址、或按条件二选一。机器指令译码在这里退化成一次跳转：取指结束后，用 opcode（必要时拼上功能位）查一张入口表，得到这条机器指令微程序的首地址，交给微序列器接着走。条件分叉也用同一机制：条件跳转指令的微程序里安排一条指定“看 Zero 标志”的微指令，Zero 为 1 时下一微地址指向“更新 PC”的微指令，为 0 时指向“直接回取指”的微指令。

\begin{lstlisting}
; 注释（推进）：每个 u/T 行代表一个控制节拍，同一行内的控制信号并行生效。
; 注释（边界）：next 决定下一微地址；等待或错误时不得提前进入写回。
; 微程序示意：ADD 与条件跳转 BEQ 各对应一段微指令
FETCH:    mem_read=1, ir_write=1, next=DISPATCH
DISPATCH: 按 opcode 查入口表 -> ADD_U0 / BEQ_U0 / ...
ADD_U0:   a_sel=REG_A, b_sel=REG_B, alu_op=ADD,
          reg_write=1, flags_write=1, next=FETCH
BEQ_U0:   a_sel=PC, b_sel=OFFSET, alu_op=ADD,
          条件选择=Zero: Z=1 -> BEQ_U1, Z=0 -> FETCH
BEQ_U1:   wb_sel=ALU_RESULT, pc_write=1, next=FETCH
\end{lstlisting}

把微程序控制器本身画成结构图，就是五个部件构成的闭环。

\begin{pdfdiagram}
  \node[hw state, minimum width=2.0cm, align=center] (uar) at (0,0) {微地址\\寄存器};
  \node[hw memory, minimum width=2.4cm, align=center] (cstore) at (3.2,0) {控制存储器\\（微程序）};
  \node[hw state, minimum width=2.0cm, align=center] (uir) at (6.5,0) {微指令\\寄存器};
  \node[hw memory, minimum width=2.4cm, align=center] (map) at (0,-2.2) {入口表\\opcode→首地址};
  \node[hw control, minimum width=2.0cm, align=center] (seq) at (3.2,-2.2) {微序列器};
  \node[hw compute, minimum width=2.2cm, align=center] (cond) at (6.5,-2.2) {条件选择\\（看标志）};
  \draw[hw bus] (uar.east) -- (cstore.west);
  \node[hw label] at (1.5,0.25) {微地址};
  \draw[hw bus] (cstore.east) -- (uir.west);
  \node[hw label] at (4.95,0.25) {微指令};
  \draw[hw bus] (uir.north) -- ++(0,0.7) node[above, hw label] {控制线字段 → 数据通路};
  \draw[hw bus] (uir.south) -- (cond.north);
  \node[hw label, anchor=west] at (6.65,-1.1) {下一/分支微地址};
  \node[hw label, align=center] (fl) at (9.0,-2.2) {标志\\Z/C/V/N};
  \draw[hw return, dashed] (fl.west) -- (cond.east);
  \draw[hw bus] (cond.west) -- (seq.east);
  \node[hw label] at (4.8,-1.95) {条件结果};
  \draw[hw bus] (map.east) -- (seq.west);
  \node[hw label] at (1.7,-1.98) {首地址};
  \node[hw label] (opc) at (0,-3.3) {opcode（取指后）};
  \draw[hw return, dashed] (opc.north) -- (map.south);
  \draw[hw bus] (seq.north) -- (3.2,-1.35) -- (-1.9,-1.35) -- (-1.9,0) -- (uar.west);
  \node[hw label] at (0.6,-1.15) {下一微地址};
\end{pdfdiagram}
\diagramnote{微程序控制器的五个部件构成闭环：微地址寄存器给出当前微地址；控制存储器读出该行的微指令；微指令寄存器的控制线字段直接驱动数据通路，下一微地址和条件选择字段交给微序列器；条件选择按标志二选一；opcode 经入口表提供各段微程序的首地址。改指令行为只改控制存储器内容，硬件不动。}

与普通组合逻辑译码（硬连线控制）对照，差别不在“控制线从哪来”——最终每根控制线都要在正确的周期输出正确的电平——而在“规则写在哪里”。硬连线方案把规则固化在与或门阵列里，改一条指令的行为就要改逻辑、重新布线；微程序方案把规则写成存储器内容，改指令行为只改微码，硬件本身一根线不动。这个差别在工程上影响深远：指令集复杂、寻址方式多的机器，比如 8086 这类变长指令、多周期动作的处理器，用硬连线译码会让组合逻辑大到难以设计和检查，微程序把复杂度转移成“填表”；机器出厂后发现指令行为有问题，有时只更新微码就能修正，不必更换芯片。代价同样明确：每条微指令多一次控制存储器访问，周期下限受存储器延迟约束；微地址、条件选择、入口表又是一层要调试的状态。所以规则简单的现代核心回到硬连线快速路径，把微码留给复杂指令和罕见情况——两层控制器并存，正是这一节两种方案在真实芯片里的最终分工。

在分立实验中，每个部件都有对应实现和各自的检查点：状态寄存器用 74HC161 计数器或 D 触发器组，检查每个时钟沿状态是否按预期推进；状态译码用 74HC138 或门阵列，检查是否只有当前状态对应的控制线有效；控制存储用 EEPROM、ROM 或拨码开关矩阵，检查同一地址是否输出稳定控制字；条件选择让 Zero/Carry/Overflow 进入下一状态逻辑，检查条件成立和不成立是否走向不同状态；复位逻辑用复位按钮加同步释放或明确初始化，检查上电后是否进入确定状态。

控制错误往往比 ALU 错误更隐蔽，因为 ALU 本身可能算对了，但错误选择线把结果送错地方，或者错误写使能在错误时钟沿提交状态。

排查要按控制线分类：\code{alu_op} 错会算错函数、标志位也可能错，先查 opcode 译码、功能位和状态选择；\code{a_sel/b_sel} 错让 ALU 输入不是预期来源，先查操作数字段、mux 控制和临时寄存器；\code{wb_sel} 错让正确结果没有被写回，先查写回 mux 和控制状态；\code{reg_write} 错让目标寄存器未写或被意外覆盖，先查写使能门控和时钟边界；\code{flags_write} 错让条件分支使用错误标志，先查标志寄存器写入条件和指令语义；\code{mem_write} 错让存储器或设备被误写，先查地址译码、写周期和保护条件。

8086 常被用来说明控制复杂度。它的指令长度不固定，寻址方式丰富，许多动作需要拆成多个内部步骤；因此不能把控制器想象成“opcode 直接连到 ALU op”这么简单。80386 进一步加入 32-bit 模式、保护机制和更复杂的地址转换相关检查，控制路径需要处理更多条件。现代处理器又在硬连线快速路径、微操作、预测、乱序调度和异常恢复之间做更复杂的工程分解。无论规模如何，控制器最终仍要回答同一个问题：本周期哪些电线有效，哪个状态会被提交。

\topic{用控制 trace 走读一次动作}

复杂动作可以拆成微操作：读、选择、计算、保存、更新状态。本节所称微操作，指一个周期内同时打开的一组硬件控制信号。以下先以 \code{R1} 和 \code{R2} 相加、结果写入 \code{R3} 为例。为了让时序更容易观察，采用三周期版本：S0 读源寄存器并锁存，S1 计算并锁存 ALU 输出，S2 写回目标寄存器。

\begin{lstlisting}
// 注释（推进）：每个 u/T 行代表一个控制节拍，同一行内的控制信号并行生效。
// 注释（边界）：next 决定下一微地址；等待或错误时不得提前进入写回。
S0: read_sel_a=R1, read_sel_b=R2
    A_write=1, B_write=1
    reg_write=0, ALU_out_write=0

S1: a_sel=A, b_sel=B, alu_op=ADD
    ALU_out_write=1
    flags_write=1
    reg_write=0, A_write=0, B_write=0

S2: wb_sel=ALU_out, write_sel=R3
    reg_write=1
    A_write=0, B_write=0, ALU_out_write=0
\end{lstlisting}

第一步 S0 需要落实为具体控制信号。\code{read_sel_a} 选中 R1 后，寄存器阵列的第一个读端口会把 R1 当前保存的 bit 组合驱动到读数据线上；\code{read_sel_b} 选中 R2 后，第二个读端口同理输出 R2。等这些组合路径稳定后，时钟沿到来，\code{A_write} 和 \code{B_write} 让 A、B 锁存器保存这两个值。这个周期里 \code{reg_write} 必须为 0，因为此时写回 mux 上的值并不代表最终结果。

第二步 S1 负责 ALU 计算。A、B 锁存器的输出作为 ALU 输入，\code{alu_op=ADD} 让 ALU 内部选择加法路径。加法电路不是瞬间得到稳定结果，低位进位会影响高位，输出线上可能经历一小段变化过程。控制器在这个周期打开 \code{ALU_out_write}，让时钟沿把已经稳定的 ALU 输出锁存下来；同时 \code{reg_write} 仍然保持 0。若本动作需要后续条件判断，\code{flags_write} 也在这个周期打开，把 Zero、Carry、Overflow 和 Negative 保存到标志寄存器。

第三步 S2 负责写回。\code{wb_sel} 选择 ALUOut 锁存器，\code{write_sel} 选择 R3，\code{reg_write} 在这个周期置 1。到时钟沿时，寄存器阵列把写数据端口上的值保存到 R3。至此，R1 和 R2 没有被改动，R3 得到 R1+R2 的结果，控制器回到空闲、取指或下一动作状态。

\begin{longtable}{L{0.12\linewidth}L{0.34\linewidth}L{0.24\linewidth}L{0.2\linewidth}}
\toprule
状态 & 打开的路径 & 周期末提交 & 必须关闭 \\
\midrule
S0 & R1/R2 到 A/B 锁存器 & A、B & \code{reg_write}、\code{ALU_out_write} \\\tablerowrule
S1 & A/B 到 ALU，ALU 到 ALUOut & ALUOut、Flags & \code{reg_write} \\\tablerowrule
S2 & ALUOut 到写回端口，写地址 R3 & R3 & A/B 写使能、ALUOut 写使能 \\
\bottomrule
\end{longtable}

把这张表画成对齐波形后，“打开/提交/关闭”的关系就能逐项检查：

\begin{waveform}[x=1cm,y=1cm]
  \draw[BookAccent, line width=0.9pt] (0,6.8) -- (2,6.8) -- (2,7.6) -- (4,7.6) -- (4,6.8) -- (6,6.8) -- (6,7.6) -- (8,7.6) -- (8,6.8) -- (10,6.8) -- (10,7.6) -- (11,7.6);
  \node[hw label, anchor=east] at (-0.2,7.2) {CLK};
  \draw[BookTeal, line width=0.9pt] (0,5.4) -- (10,5.4) -- (10,6.2) -- (0,6.2) -- cycle;
  \foreach \x in {2,4,6,8} { \draw[BookTeal, line width=0.9pt] (\x,5.4) -- (\x,6.2); }
  \node at (1.0,5.8) {S0};
  \node at (3.0,5.8) {S1};
  \node at (5.0,5.8) {S2};
  \node at (7.0,5.8) {下一动作};
  \node[hw label, anchor=east] at (-0.2,5.8) {状态};
  \draw[BookAccent, line width=0.9pt] (0,4.2) -- (2,4.2) -- (2,5.0) -- (4,5.0) -- (4,4.2) -- (11,4.2);
  \node[hw label, anchor=east] at (-0.2,4.6) {\scriptsize A\_write, B\_write};
  \draw[BookAccent, line width=0.9pt] (0,3.0) -- (4,3.0) -- (4,3.8) -- (6,3.8) -- (6,3.0) -- (11,3.0);
  \node[hw label, anchor=east] at (-0.2,3.4) {\scriptsize ALU\_out\_write, flags\_write};
  \draw[BookAccent, line width=0.9pt] (0,1.8) -- (6,1.8) -- (6,2.6) -- (8,2.6) -- (8,1.8) -- (11,1.8);
  \node[hw label, anchor=east] at (-0.2,2.2) {\scriptsize reg\_write};
  \draw[BookMuted, line width=0.9pt] (0,0.4) -- (10,0.4) -- (10,1.2) -- (0,1.2) -- cycle;
  \foreach \x/\v in {2.3/{R1, R2}, 4.3/{R1, R2}, 6.3/{R1, R2}, 8.3/{R1, R2}} {
    \draw[BookMuted, line width=0.9pt] (\x,0.4) -- (\x,1.2);
    \node at ({\x+0.85},0.8) {\v};
  }
  \node at (1.15,0.8) {旧值};
  \node[hw label, anchor=east] at (-0.2,0.8) {A, B};
  \draw[BookMuted, line width=0.9pt] (0,-1.0) -- (10,-1.0) -- (10,-0.2) -- (0,-0.2) -- cycle;
  \draw[BookMuted, line width=0.9pt] (4.3,-1.0) -- (4.3,-0.2);
  \node at (2.15,-0.6) {旧值};
  \node at (7.15,-0.6) {R1+R2};
  \node[hw label, anchor=east] at (-0.2,-0.6) {ALUOut};
  \draw[BookMuted, line width=0.9pt] (0,-2.4) -- (10,-2.4) -- (10,-1.6) -- (0,-1.6) -- cycle;
  \draw[BookMuted, line width=0.9pt] (6.3,-2.4) -- (6.3,-1.6);
  \node at (3.15,-2.0) {旧值};
  \node at (8.15,-2.0) {R1+R2};
  \node[hw label, anchor=east] at (-0.2,-2.0) {R3};
  \foreach \x in {2,4,6,8} { \draw[BookMuted, dashed, line width=0.4pt] (\x,-2.7) -- (\x,7.9); }
\end{waveform}
\diagramnote{三周期 \code{R3=R1+R2} 的对齐波形：控制线只在规定周期为高；状态元件只在上升沿提交。注意三条控制线的互斥关系——任何一拍 \code{reg\_write} 提前变高，R3 都会被未成熟的值污染；任何一拍 \code{ALU\_out\_write} 忘记关闭，ALUOut 锁存器就会混入旧值。}

该例也说明控制错误的风险。若 S2 的 \code{wb_sel} 错选成 B 锁存器，R3 写入的就是 R2，而不是相加结果；若 S1 提前打开 \code{reg_write}，R3 可能在 ALU 输出尚未稳定时被覆盖；若 S0 忘记关闭 \code{ALU_out_write}，旧的计算结果会被无关读数污染。微操作表的价值在于：每一步都能检查哪些线打开、哪些线关闭、哪个状态元件会在时钟沿改变。

再看一个条件动作：若 \code{R1-R2} 的结果为 0，则下一状态跳到 \code{TARGET}，否则继续 \code{NEXT}。硬件上这仍然是数据通路和控制器配合：ALU 执行 SUB，Zero 标志保存比较结果，下一状态逻辑读取 Zero 并选择目标。

\begin{lstlisting}
// 注释（推进）：每个 u/T 行代表一个控制节拍，同一行内的控制信号并行生效。
// 注释（边界）：next 决定下一微地址；等待或错误时不得提前进入写回。
S0: read_sel_a=R1, read_sel_b=R2
    A_write=1, B_write=1

S1: a_sel=A, b_sel=B, alu_op=SUB
    flags_write=1

S2: if Zero then next_state=TARGET
    else next_state=NEXT
    reg_write=0
\end{lstlisting}

这个 trace 说明比较不是神秘动作。相等比较可以由减法和 Zero 标志实现；条件跳转不是 ALU 输出直接“跳走”，而是控制器根据标志选择下一状态或下一 PC。

\topic{把算术与数据通路相关章节接到最小 CPU}

到这里，CPU 的几个必要部件已经在概念上齐备。寄存器堆保存通用值，ALU 产生算术逻辑结果和标志位，数据通路用 mux 和总线把源、结果和写回连起来，控制器在每个周期输出一组控制线。最小 CPU 还需要 PC、指令寄存器、指令译码、存储器接口和复位入口。PC 是带加法和选择的寄存器，需要顺序取指、分支目标和异常入口；指令寄存器是保存当前指令 bit 的状态元件，需要字段划分和 opcode 译码；译码器是把 bit 字段翻译成控制线的组合逻辑；存储器接口是另一条受控数据通路，需要处理总线等待、对齐、MMIO 和异常。

把这些部件真正拼成最小 CPU 时，接口清单比概念名称更重要：每个部件不只要有名称，还要有输入、输出和写入边界。

\begin{longtable}{L{0.18\linewidth}L{0.34\linewidth}L{0.38\linewidth}}
\toprule
部件 & 关键输入/输出 & 必要控制线 \\
\midrule
PC & 输出取指地址，输入顺序 PC 或分支目标 & \code{pc_write}、\code{pc_sel} \\\tablerowrule
指令寄存器 IR & 输入存储器返回的指令 bit，输出 opcode 和寄存器字段 & \code{ir_write} \\\tablerowrule
寄存器堆 & 输出两个源操作数，输入写回数据 & \code{read_addr_a/b}、\code{write_addr}、\code{reg_write} \\\tablerowrule
立即数扩展器 & 输入指令字段，输出零扩展或符号扩展值 & \code{imm_kind} 或译码规则 \\\tablerowrule
ALU & 输入两个操作数，输出结果和标志 & \code{a_sel}、\code{b_sel}、\code{alu_op}、\code{flags_write} \\\tablerowrule
存储器接口 & 地址、写数据、读数据、等待信号 & \code{mem_read}、\code{mem_write}、\code{addr_sel} \\\tablerowrule
控制器 & 输入 opcode、状态、标志和等待信号，输出控制字 & \code{next_sel}、复位入口、异常入口 \\
\bottomrule
\end{longtable}

若这张接口表无法填写完整，最小 CPU 就还只是部件堆叠，而不是可执行机器。特别要注意两条线：第一，PC、IR、寄存器堆和标志寄存器都是状态元件，必须有明确写使能；第二，存储器接口可能比 CPU 内部组合路径慢，因此控制器必须能等待、重试或保持状态。

调试整机硬件时，不要先问“CPU 为什么不工作”，而要按 trace 缩小范围：PC 是否在时钟沿更新，指令寄存器是否保存了正确 bit，译码是否输出了预期控制字，ALU 输入是否来自正确源，ALU 输出和标志是否正确，写回 mux 是否选对，写使能是否只在目标周期打开。

\topic{从机器级表达式读回硬件动作}
高级语言里的 \code{x+y}、\code{a[i]}、\code{p->field}、\code{x \& mask} 看起来属于软件语法，但映射到硬件上，它们都要变成寄存器读、立即数扩展、移位、加法、逻辑运算、LOAD/STORE、标志位和写回。

把常见表达式逐个拆开：\code{x+y} 读两个寄存器，ALU 加法，结果写回，对应 \code{alu\_op=ADD}，并产生 Carry、Overflow、Zero 标志；\code{x-y} 让 B 取反加一、复用加法器，\code{sub} 同时控制 XOR 和最低位进位；\code{x<<k} 由移位器按移位量选择各位来源，逻辑移位、算术移位和溢出规则要分开；\code{x \& mask} 走 ALU 逻辑路径保留部分 bit，常用于字段提取、权限位和状态位；\code{a[i]} 由基址加缩放下标形成地址再 LOAD/STORE，经过移位器或乘法路径、地址加法和字节使能；\code{p->field} 是指针基址加固定字段偏移，依赖结构体布局、对齐和访问宽度。

以 \code{a[i] = a[i] + 1} 为例，若 \code{a} 是 32-bit 整数数组，机器动作至少包含五步。第一步下标缩放：\code{i<<2} 或地址生成单元计算 \code{i*4}，元素大小必须由类型布局决定。第二步地址形成：基址与缩放下标相加，要说明加法按地址宽度截断还是产生异常。第三步读取旧值：发起 32-bit LOAD，字节序、对齐、cache/MMIO 属性都要清楚。第四步加一：ALU 执行 \code{old+1}，要说明是否更新条件码、溢出如何解释。第五步写回内存：发起 32-bit STORE，字节使能必须正确，不能破坏相邻元素。每一步都能追到本章已有部件。

\begin{lstlisting}
// 注释（推进）：按行追踪数据来源、经过的部件和最终写入目标。
// 注释（边界）：只有标出的时钟沿、阶段或异常入口会改变体系结构可见状态。
addr = base(a) + i * 4
old  = MEM[addr]
new  = old + 1
MEM[addr] = new
\end{lstlisting}

如果把 \code{a[i]} 错看成“数组对象知道自己的长度”，就会误解越界访问。硬件只看到基址、下标、缩放和加法。除非后续保护机制或编译器插入检查，否则 \code{i} 超出范围仍然会形成一个地址，并继续发起 LOAD 或 STORE。这个地址可能指向同一栈帧中的另一个变量、堆对象的元数据、返回地址、MMIO 区域，或未映射页。异常和分页机制正是把“地址是否允许访问”提升为硬件可检查的约定。

\topic{结构体布局把字段名变成偏移}
结构体字段名不是硬件对象。编译后，字段访问会变成“对象起始地址加字段偏移”。字段偏移由字段顺序、字段宽度和对齐规则决定。对齐不是排版习惯，而是为了让硬件访问落在更自然的字节边界上，减少拆分传输、简化字节使能和提高 cache 行利用率。

\begin{lstlisting}
struct Packet {
  uint8_t  type;
  uint8_t  flags;
  uint16_t length;
  uint32_t address;
};

p->address  ->  LOAD32 [base(p) + 4]
\end{lstlisting}

\begin{longtable}{L{0.16\linewidth}L{0.2\linewidth}L{0.24\linewidth}L{0.3\linewidth}}
\toprule
字段 & 宽度 & 可能偏移 & 硬件访问 \\
\midrule
\code{type} & 1 字节 & 0 & 字节 LOAD，通常零扩展 \\\tablerowrule
\code{flags} & 1 字节 & 1 & 字节 LOAD，可与掩码配合 \\\tablerowrule
\code{length} & 2 字节 & 2 & 半字 LOAD，受字节序影响 \\\tablerowrule
\code{address} & 4 字节 & 4 & 字 LOAD，要求字节使能覆盖四个车道 \\
\bottomrule
\end{longtable}

若为了节省空间强行取消对齐，某些字段可能跨越自然边界。硬件仍可通过多个字节车道或多个总线周期完成访问，但代价会增加；某些架构会直接规定未对齐访问异常，要求软件拆成多个字节访问。教学阶段要把这件事写清楚：类型布局不是只影响内存大小，它会进入地址形成、字节选择、cache line 命中和总线传输。

\topic{位掩码是最小的机器级抽象工具}
许多硬件状态寄存器不是把每个状态放进独立地址，而是把多个 1-bit 字段压在同一个字中。程序读出一个状态寄存器后，常用掩码提取某些位；写控制寄存器时，也常用掩码只改指定字段。掩码操作本质上是 ALU 的 AND、OR、XOR 和移位路径。

\begin{lstlisting}
// 注释（推进）：逐行追踪发起者、所有权和可见性；请求被接受不等于事务完成。
// 注释（边界）：以采样沿、状态位或 completion 为准，再允许消费者使用结果。
status = MMIO[UART_STAT]
ready  = (status & 0x01) != 0
error  = (status & 0x0E) >> 1

ctrl = ctrl | 0x04      // set bit 2
ctrl = ctrl & ~0x04     // clear bit 2
ctrl = ctrl ^ 0x04      // toggle bit 2
\end{lstlisting}

四种位操作各走一条 ALU 路径：AND mask 保留掩码为 1 的位、清除其他位，用于提取状态位、对齐地址和清除字段；OR mask 把掩码为 1 的位置 1，用于设置控制位和构造权限字段；XOR mask 翻转掩码为 1 的位，用于切换状态和构造取反路径；SHIFT 改变字段位置或完成缩放，用于字段对齐和数组下标乘元素宽度。

位掩码也能解释为什么设备寄存器规格必须写清楚读写语义。有些位是只读状态位，有些位写 1 清除，有些位写 0 无效，有些位一写就启动设备动作。若用普通“读-改-写”序列处理一个带副作用寄存器，可能在读时清掉事件，或在写回时误改其他字段。同样是 bit，普通 RAM 字节和设备寄存器字段的约定完全不同。

\topic{整数溢出要同时给出数学结果和机器结果}
固定宽度整数运算的结果总是某个位宽内的 bit 模式。数学上的结果可能超出范围，此时硬件仍然给出低位截断结果和若干标志位。若程序把 \code{INT_MAX+1} 当作普通正数继续使用，后续比较、数组下标和地址形成都会被污染。硬件不会自动替程序选择“报错、饱和、回绕、陷入”哪一种语义，必须由 ISA、编译器、语言规则或显式指令规定。

若 ISA 提供普通 ADD、带陷入 ADD、饱和 ADD 或宽乘法指令，控制器和 ALU 就必须输出不同的控制线和标志处理路径。若 ISA 只提供普通 ADD，软件就需要用比较和条件分支组合出额外检查。所谓“语言整数语义”，最后仍要体现为硬件结果、标志位、异常入口和控制流。

四种溢出语义各适合一类场景：回绕只保留低 N 位并由标志位记录越界，适合地址计数器、哈希和无符号模运算；陷入在溢出时进入异常或错误路径，适合安全语言运行时和调试检查；饱和把超范围结果钳到最大或最小值，适合 DSP、音频、图像和控制系统；扩大位宽用更宽临时结果保存数学值，适合乘法、累加和编译器优化的中间值。

\topic{案例：解析一个二进制包头}
系统程序经常要解析来自网络、磁盘或外设 FIFO（先进先出缓冲）的二进制包头。这个任务看起来属于软件格式处理，实际上集中使用了算术与数据通路相关章节的全部机器级工具：字节序、字段偏移、位掩码、整数扩展、对齐访问和溢出检查。假设某设备返回如下 8 字节包头：

\begin{lstlisting}
offset  size  field
0       1     type
1       1     flags
2       2     length   // little-endian
4       4     address  // little-endian
\end{lstlisting}

解析过程可以写成：

\begin{lstlisting}
type   = LOAD8 [base + 0]
flags  = LOAD8 [base + 1]
len_lo = LOAD8 [base + 2]
len_hi = LOAD8 [base + 3]
length = len_lo | (len_hi << 8)
addr   = LOAD32 [base + 4]
ready  = (flags & 0x01) != 0
error  = (flags & 0x0E) >> 1
\end{lstlisting}

\begin{longtable}{L{0.18\linewidth}L{0.38\linewidth}L{0.34\linewidth}}
\toprule
解析动作 & 硬件路径 & 需要确认的问题 \\
\midrule
读取 \code{type} & 字节 LOAD 后零扩展 & 目标寄存器高位是否清零 \\\tablerowrule
读取 \code{length} & 两个字节按小端拼接 & 字节序是否和设备规格一致 \\\tablerowrule
读取 \code{address} & 32-bit LOAD 或四次字节 LOAD 组合 & 对齐、字节使能和异常规则是否清楚 \\\tablerowrule
提取 \code{ready} & AND 掩码后比较零 & 状态位是否被读操作清除 \\\tablerowrule
提取 \code{error} & AND 后右移 & 字段位置和移位方向是否正确 \\\tablerowrule
检查长度 & 与缓冲区容量比较 & 无符号比较和溢出检查不能混用 \\
\bottomrule
\end{longtable}

若 \code{length} 来自外部设备，不能直接用它做 \code{base + length}。地址加法可能回绕，长度可能超过缓冲区，字段可能被设备错误写入。可靠代码需要先做范围检查，再执行后续 LOAD/STORE。硬件层面，这些检查仍然是比较器、标志位和分支。

\begin{lstlisting}
// 注释（推进）：先判断条件，再只执行命中的分支；分支顺序同时表达优先级。
// 注释（边界）：所有输入稳定后才读取结果；未覆盖的条件必须有明确默认行为。
if length > buffer_capacity:
  reject packet
if base + length < base:
  reject packet      // unsigned wraparound check
payload_end = base + length
\end{lstlisting}

\begin{longtable}{L{0.2\linewidth}L{0.36\linewidth}L{0.34\linewidth}}
\toprule
检查项 & 机器级机制 & 漏掉后的后果 \\
\midrule
长度上限 & 无符号比较 \code{length <= capacity} & 后续 STORE 越过缓冲区 \\\tablerowrule
地址回绕 & \code{base+length < base} 说明无符号回绕 & 访问低地址或错误对象 \\\tablerowrule
字段合法性 & 掩码后剩余位应为 0 或规定值 & 未定义设备状态被当成合法状态 \\\tablerowrule
对齐要求 & \code{address \& (align-1)==0} & 未对齐异常或拆分传输代价 \\\tablerowrule
权限属性 & 地址不能落入 MMIO 或只读区域 & 普通数据访问变成设备副作用或保护异常 \\
\bottomrule
\end{longtable}

每一步都要指出需要哪些 LOAD、移位、掩码、比较、标志和分支；这些动作失败时，会表现为总线错误、DMA 可见性问题，或页错误和保护异常。

\topic{案例：多字长整数为什么要保存 Carry}
很多低成本 CPU 或教学 CPU 的 ALU 宽度较小，但程序需要处理更宽整数。例如用 16-bit ALU 实现 32-bit 加法，必须先加低 16 位，再把低位加法产生的 Carry 加入高 16 位。这里 Carry 不是附属显示灯，而是下一步运算的真实输入。

\begin{lstlisting}
// 32-bit A+B on a 16-bit ALU
lo = A_lo + B_lo
carry = carry_out(lo)
hi = A_hi + B_hi + carry
result = [hi:lo]
\end{lstlisting}

减法、多精度计数器、地址加法和校验和也会用到类似链路。若控制器在低位加法后错误关闭 \code{flags_write}，或在高位加法前让另一条指令覆盖 Carry，多字长结果就会错。真实 ISA 提供 ADC、SBB 这类“带进位加/减”指令，正是因为标志位是程序员可见状态的一部分。

把这条链路的控制线逐拍列清：低位加法用 \code{alu\_op=ADD}、\code{carry\_in=0}，预期低位结果和 CarryOut 正确；保存 Carry 用 \code{flags\_write=1}，预期标志寄存器在下一周期仍可读；高位带进位加法用 \code{carry\_in=Carry}，预期高位结果包含低位进位；写回组合结果时写回低位和高位目标寄存器，预期两个半字都在正确的时钟边界提交。

\topic{案例：32-bit 加法拆进 8-bit ALU}
上一个案例用 16-bit ALU 分两次完成 32-bit 加法。若数据通路只有 8-bit 宽——早期微处理器和许多单片机正是如此——同一个 32-bit 加法就要拆成四步，从最低字节开始逐字节推进。每步的差别只有两处：第一步的最低字节没有进位输入，用普通 ADD；后面三步都必须把上一步的 Carry 当作本次加法的第三个输入，用带进位加 ADC。拆字长不是把同一个动作简单重复四遍，而是一条穿过四个字节的进位接力。

以 $A=\code{0x1FF34A80}$、$B=\code{0x0E21C590}$ 为例，逐字节 trace 如下：

\begin{longtable}{L{0.1\linewidth}L{0.34\linewidth}L{0.16\linewidth}L{0.3\linewidth}}
\toprule
步骤 & 本字节计算 & 结果字节 & Carry 输出与去向 \\
\midrule
1 & \code{0x80+0x90}（ADD，进位输入为 0） & \code{0x10} & Carry=1，送入步骤 2 \\\tablerowrule
2 & \code{0x4A+0xC5+1}（ADC，加上一步 Carry） & \code{0x10} & Carry=1，送入步骤 3 \\\tablerowrule
3 & \code{0xF3+0x21+1}（ADC） & \code{0x15} & Carry=1，送入步骤 4 \\\tablerowrule
4 & \code{0x1F+0x0E+1}（ADC） & \code{0x2E} & Carry=0，全和无第 33 位 \\
\bottomrule
\end{longtable}

四步之后从高到低拼出 \code{0x2E151010}。逐步核对数值：步骤 1 的 $128+144=272$，结果字节只保留 \code{0x10}，超出的 256 变成 Carry；步骤 2 的 $74+197+1=272$，同样保留 \code{0x10} 并把 1 交给下一步；步骤 3 的 $243+33+1=277$，保留 \code{0x15}；步骤 4 的 $31+14+1=46$，恰好不再越界。若步骤 3 漏掉来自步骤 2 的 Carry，本字节就差 1，错误还会沿进位链继续向高字节放大。这条链路上的每个 Carry 都不是显示信息，而是下一步的真实数据输入：任何一步算完后 Carry 被意外覆盖，后面所有字节的和都不可信。

这就是 ADC 指令存在的理由。普通 ADD 的进位输入恒为 0，无法表达“把上一步的进位加进来”；ISA 若只提供 ADD，软件就要用额外的比较和条件分支手工重建进位，代价远高于一条直接读取 Carry 标志的加法指令。因此从 8080、8086 到 x86-64，指令集都把 ADD 与 ADC 成对提供：做超出寄存器宽度的加法时，最低位用 \code{add}，更高位依次用 \code{adc}，且两条指令之间不能插入会改写标志的第三条指令。硬件层面，ADD 与 ADC 往往共享同一条进位链，差别只在 \code{carry\_in} 这一根线接 0 还是接标志寄存器的 Carry 输出：

\begin{lstlisting}
byte0: alu_op=ADD, carry_in=0,   flags_write=1
byte1: alu_op=ADC, carry_in=CF,  flags_write=1
byte2: alu_op=ADC, carry_in=CF,  flags_write=1
byte3: alu_op=ADC, carry_in=CF,  flags_write=1
\end{lstlisting}

宽数据被切成机器字长的若干片，片与片之间靠标志寄存器里的一位 Carry 衔接。多字长运算的正确性，最终就取决于这一位能否在两次 ALU 动作之间原样存活。

\topic{案例：把十进制字符串转成二进制整数}
程序从串口、文件或命令行读到的数字是一串字符，不是数值。字符 \code{'0'} 到 \code{'9'} 在 ASCII 中的编码是 \code{0x30} 到 \code{0x39}，连续排列，所以“字符减 \code{0x30}”就得到该数字的数值 0 到 9。剩下的事情是按十进制位权组装：每读到一个新数字，已累积的结果扩大十倍，再加上新数值，即 \code{result = result*10 + digit}。以字符串 \code{"1234"} 为例：

\begin{longtable}{L{0.18\linewidth}L{0.26\linewidth}L{0.46\linewidth}}
\toprule
读到字符 & 减 \code{0x30} 得数值 & 累积结果 \code{result*10+digit} \\
\midrule
\code{'1'} = \code{0x31} & 1 & $0\times10+1=1$ \\\tablerowrule
\code{'2'} = \code{0x32} & 2 & $1\times10+2=12$ \\\tablerowrule
\code{'3'} = \code{0x33} & 3 & $12\times10+3=123$ \\\tablerowrule
\code{'4'} = \code{0x34} & 4 & $123\times10+4=1234$ \\
\bottomrule
\end{longtable}

四步之后得到二进制整数 $1234=\code{0x04D2}$。循环里唯一的算术是“乘十加”，而乘十不必动用乘法器：$10x=8x+2x$，乘 8 是左移 3 位，乘 2 是左移 1 位，两次移位加两次加法即可。移位器和加法器是 ALU 里已有的路径，乘法器则是面积大、延迟长的独立部件；编译器把常数乘法改写成移位加组合，正是强度削减优化的典型例子。对应的 x86-64 汇编（Intel 语法）如下，设 \code{rsi} 指向字符串、遇到第一个非数字字符结束：

\begin{lstlisting}
    xor  eax, eax          ; result = 0
.next:
    movzx ecx, byte [rsi]  ; 读当前字符并零扩展
    sub   ecx, 0x30        ; '0' 的编码是 0x30，得数值 0..9
    cmp   ecx, 9
    ja    .done            ; 无符号大于 9：非数字，结束
    mov   rdx, rax
    shl   rdx, 3           ; rdx = result * 8
    shl   rax, 1           ; rax = result * 2
    add   rax, rdx         ; rax = result * 10
    add   rax, rcx         ; rax = result * 10 + digit
    inc   rsi
    jmp   .next
.done:
\end{lstlisting}

这段代码值得用本章的眼光再读一遍。\code{sub} 与 \code{cmp} 共用 ALU 的减法路径并更新标志；\code{ja} 是无符号条件跳转，依据 Carry 标志以及 Zero 标志（CF=0 且 ZF=0 才跳转）——字符小于 \code{'0'} 时减法产生借位、结果成为很大的无符号数，与大于 9 的情形被同一次比较一并拦下，所以一次无符号比较就完成了上下两个边界的检查。\code{shl} 走桶形移位器路径，\code{add} 走进位链，整个循环没有任何一条指令需要乘法器。

\topic{专题：从门级代价看 ALU 为什么这样组织}
ALU 表面上像一个会做很多运算的黑盒，内部通常更像若干条专用组合路径加一个结果选择器。加法/减法依赖进位链，逻辑运算依赖逐位门，移位依赖多级选择网络，比较往往复用减法路径和标志。\code{ADD}、\code{AND}、\code{SLL}、\code{SLT} 走的是不同硬件路径，代价并不相同。

各项功能的来源和代价也不同：加法来自全加器链、超前进位或前缀进位网络，进位传播决定关键路径，位宽越大越明显；减法让 \code{B} 逐位取反并令最低位 \code{carry\_in=1}，Carry 和 Overflow 的解释必须按减法规则重新定义；按位逻辑是每一位独立的 AND/OR/XOR/NOT，延迟短，但结果选择 mux 会增加负载；无符号小于看减法后的借位或 Carry 标志，不能直接看符号位；有符号小于由减法结果符号与 Overflow 共同决定，\code{0x8000-1} 这类边界最容易出错；移位用桶形移位器或多周期逐位移位，单周期桶形移位器面积和 mux 级数较大。

如果教学 CPU 的目标是先跑通控制 trace，可以先用串行进位（纹波进位）加法器和多周期移位；如果目标是提高主频，就要把加法器换成 carry-lookahead、carry-select 或 prefix 结构。这里的选择没有绝对优劣，只有约束不同：面积、功耗、最高频率、验证复杂度和 ISA 对单周期完成的承诺不同。

进位优化的核心是把“等待前一位进位”变成“并行计算哪些位会产生或传播进位”，所用的就是前面两节反复出现的 \(G\)、\(P\) 信号和同一条进位递推 \(C_{i+1}=G_i\lor(P_iC_i)\)。前缀加法器会把这些关系按树形组合，缩短最长逻辑深度，但布线和门数会上升。一个 64-bit ALU 不能只按“64 个全加器排队”等比例粗略估计性能。

\topic{专题：ALU 状态标志的生成电路}
Zero、Carry、Overflow、Negative 四个标志容易被初学者想象成“算完结果之后，再拿结果去分析一遍”的后置步骤。硬件上恰恰相反：四个标志都是同一次运算的副产物，与结果共用同一批内部信号，在同一个周期内与结果一起稳定。它们不是事后的评论，而是进位链和结果线上早已存在的信息，各自由一小段专用组合电路引出。

四个标志的来源各不相同。Negative 最简单，就是结果最高位那根线本身，补码解释下它天然指示正负，不经过任何门。Zero 是结果全部位经 OR 树归并后再取反：任何一位为 1 都会让 OR 树输出 1，只有结果全 0 时 Zero 才为 1，其树深随位宽按对数增长。Carry 来自进位链的最后一级——加法时是最高位向外的进位，减法时按约定改读为借位；它通常是标志中稳定最晚的一个，因为进位链本来就是 ALU 的关键路径。Overflow 是符号位两处进位的异或：记最高位的进位输入为 \(C_{n-1}\)、进位输出（即 Carry）为 \(C_n\)，则 \(V=C_{n-1}\oplus C_n\)——两个正数相加时，若数值部分向符号位进了位而符号位没有向外进位，符号就被异常翻转，异或正好检出这种“两处进位不一致”。

\begin{longtable}{L{0.16\linewidth}L{0.44\linewidth}L{0.3\linewidth}}
\toprule
标志 & 生成电路 & 延迟特点 \\
\midrule
Negative & 结果最高位直接引出，不经过任何门 & 与结果线同时稳定 \\\tablerowrule
Zero & 全部结果位的 OR 树加一级取反 & 树深随位宽对数增长 \\\tablerowrule
Carry & 进位链末级的进位输出 & 与进位链同长，常是关键路径 \\\tablerowrule
Overflow & 最高位进位输入与进位输出相异或 & 在 Carry 之后只多一级 XOR \\
\bottomrule
\end{longtable}

理解“同周期产生”对控制器设计有直接影响。前面三周期 trace 的 S1 拍里，\code{flags\_write} 与 \code{ALU\_out\_write} 在同一个时钟沿把标志和结果分别锁存进标志寄存器和 ALUOut，二者谁也不等谁；S2 拍的条件分支读到的标志，与写回的数据严格属于同一次运算。若把标志想象成下一周期才慢慢算出来的东西，就会错误地在控制器里多加等待拍，或者反过来在标志尚未稳定的周期提前采样。真实时序里要照顾的只是各标志延迟不同：Carry 和 Overflow 沿进位链走，稳定最晚，标志寄存器的建立时间预算必须按它们来定。

\topic{专题：乘法、除法和规格化为什么常常多周期}
一个 \(n\) 位乘法可以看成 \(n\) 个部分积的移位求和；一个除法可以看成反复试减、恢复或不恢复余数；浮点运算还要经历对阶、尾数运算、规格化和舍入。若所有步骤都做成单周期组合逻辑，面积和延迟会迅速失控。

每种运算的拆分方式和控制方式也不同：无符号乘法判断乘数最低位、条件加被乘数、移位、计数，用多周期移位加法或阵列/树形乘法器实现；补码乘法要处理符号扩展和负权部分积，Booth 编码能减少连续 1 的部分积数量；无符号除法左移余数、试减除数、决定商位、必要时恢复余数，用多周期状态机逐位产生商；定点乘法在原始整数相乘后调整缩放位数，乘积位宽翻倍，截断和舍入必须规定；浮点加法要对阶、尾数加减、规格化、舍入、异常标志，分流水级或由微码/协处理器完成。

以 8-bit 乘法为例，若采用移位加法，硬件只需要一个加法器、若干寄存器和一个计数器。每拍检查乘数最低位：若为 1，把当前被乘数加到累加器；然后被乘数左移，乘数右移。8 拍后得到 16-bit 乘积。这个实现慢，但面积小，特别适合早期处理器或低成本控制器。

\begin{lstlisting}
// 注释（推进）：每轮只推进一个元素或一个控制步，并保持中间状态的一致性。
// 注释（边界）：退出条件成立后才提交结果；空集合、越界和失败要单独处理。
product = 0
for i in 0..7:
  if (multiplier[0] == 1):
    product = product + multiplicand
  multiplicand = multiplicand << 1
  multiplier = multiplier >> 1
\end{lstlisting}

除法值得给一个完整例子。以 4-bit 恢复余数除法计算 $7 \div 2$：余数寄存器 R 初值为 0，被除数寄存器 Q 初值为 \code{0111}（7）。每一拍先把 (R,Q) 整体左移一位，再试减除数；试减结果为负说明“这一位商不上”，商位记 0 并把余数加回去（恢复），否则商位记 1：

\begin{longtable}{L{0.12\linewidth}L{0.22\linewidth}L{0.24\linewidth}L{0.14\linewidth}L{0.18\linewidth}}
\toprule
拍 & (R,Q) 左移后 & 试减 \code{R-2} & 商位 & 动作 \\
\midrule
1 & \code{0000,1110} & $0-2<0$ & 0 & 恢复 R=\code{0000} \\\tablerowrule
2 & \code{0001,1100} & $1-2<0$ & 0 & 恢复 R=\code{0001} \\\tablerowrule
3 & \code{0011,1000} & $3-2=1$ & 1 & 保留 R=\code{0001} \\\tablerowrule
4 & \code{0011,0010} & $3-2=1$ & 1 & 保留 R=\code{0001} \\
\bottomrule
\end{longtable}

四拍后 Q=\code{0011}=3，R=\code{0001}=1：$7 \div 2 = 3$ 余 1。每一步只用移位、减法和一次条件恢复——这就是为什么除法天然是多周期状态机：每拍都有“移位、试减、判负、恢复”四个微操作，商逐位产生，没有一条能把所有位一次算完的组合捷径。

浮点的难点更隐蔽。两个指数不同的浮点数相加时，较小数的尾数要右移对齐；右移时丢掉的位会进入 guard、round、sticky 三个附加位——guard 是被移出的最高一位，round 是次一位，sticky 是其余被移出位按 OR 合并出的“是否还有 1”，舍入电路用这三个附加位决定末位该不该进 1。尾数相加后可能需要左规或右规；最后按舍入模式提交。浮点不是“整数加法多一个小数点”：对阶、舍入和异常标志解释了为什么 \((a+b)+c\) 和 \(a+(b+c)\) 可能不同，也解释了 NaN、Inf、非规格化数的存在。

\topic{专题：规格化左移右移的代价}
浮点加减法把尾数对齐并相加之后，工作还没有结束：和或差不一定仍保持规格化形式 \(1.f\times2^{e}\)。尾数相加可能产生向整数位的进位，例如 \(1.110+1.011=11.001\)，结果大于等于 2；尾数相减可能让前导 1 消失，例如 \(1.001-1.000=0.001\)，整数位变成 0。两种情况都必须把结果移回“小数点前恰好一个 1”的形式，并同步调整指数，这一步叫规格化。前一种情形右移一位、指数加一，称为右规；后一种情形左移若干位、指数减去相应位数，称为左规。

两种规格化对应两条不同的硬件路径。右规的情形简单：进位只可能多出一位，硬件只需要一条固定右移一位的通路，把多出的进位放回最高位、把末位挤进 guard 位，同时给指数加一。左规的情形麻烦得多：相减造成的抵消可能抹掉任意多个前导位，最坏情况下（两个几乎相等的数相减）需要左移几十位。因此左规路径必须先在尾数上做前导零计数——用优先级编码器找出第一个 1 的位置 \(k\)——再按 \(k\) 左移，并从指数中减去 \(k\)。这条“先探测移位量、再移位”的链是浮点加法器特有的组合逻辑，定点 ALU 里并不存在。

把这两条路径画出来，左规比右规多出的那级前导零计数就清楚了。

\begin{pdfdiagram}
  \node[hw state, minimum width=2.6cm, align=center] (sum) at (0,0) {尾数相加结果\\（含进位）};
  \node[hw compute, minimum width=3.0cm, align=center] (lzc) at (3.3,1.3) {前导零计数\\（优先级编码器）};
  \node[hw compute, minimum width=2.9cm, align=center] (bs) at (7.0,0) {桶形移位器\\左规 $k$ / 右规 1\\五层 mux};
  \node[hw compute, minimum width=2.6cm, align=center] (adj) at (3.3,-2.2) {指数调整\\$e+1$ 或 $e-k$};
  \draw[hw bus] (sum.east) -- (bs.west);
  \node[hw label] at (2.2,0.18) {未规格化尾数};
  \draw[hw bus] (sum.north) -- (0,1.3) -- (lzc.west);
  \node[hw label] at (0.9,1.48) {尾数};
  \draw[hw bus] (lzc.east) -- (7.0,1.3) -- (bs.north);
  \node[hw label] at (5.9,1.48) {移位量 $k$};
  \draw[hw bus] (lzc.south) -- (adj.north);
  \node[hw label, anchor=west] at (3.45,-0.8) {$k$};
  \draw[hw bus] (bs.east) -- ++(0.8,0) node[above, hw label] {尾数 $1.f$};
  \node[hw label] (e0) at (0.6,-2.2) {原指数 $e$};
  \draw[hw bus] (e0.east) -- (adj.west);
  \draw[hw bus] (adj.east) -- ++(0.8,0) node[right, hw label] {新指数};
\end{pdfdiagram}
\diagramnote{规格化的两条硬件路径：尾数和送入桶形移位器的同时，也送入前导零计数（优先级编码器）找出第一个 1 的位置 $k$；左规按 $k$ 左移、指数减 $k$，右规只需固定右移一位、指数加一。单精度 24 位尾数用移 1、2、4、8、16 五层 mux，任意 0 到 23 位左移在一轮组合逻辑内完成。}

移位量不固定正是需要桶形移位器的原因。若用逐位移位寄存器，左规最坏要花费与尾数位宽相等的周期数，单精度是 24 位、双精度是 53 位，而且延迟随数据变化，流水级无法按固定节拍安排。桶形移位器用 \(\log_{2}n\) 层 mux 按二进制分解移位量：24 位尾数用移 1、2、4、8、16 五层，每层由移位量的一个 bit 控制，任意 0 到 23 位的左移都在一轮组合逻辑内完成，延迟固定且可预测。代价是面积和布线：每层都要把每个输出位接到多个候选输入之一，\(n\) 位 \(\log_{2}n\) 层共需约 \(n\log_{2}n\) 个二选一开关；对阶右移和左规两条路径还各占一套。

规格化还反过来影响舍入。右规把一位移入 guard 位，左规把 guard、round、sticky 里的信息带回尾数低位，因此对阶时算好的舍入依据在规格化之后必须重新评估，舍入电路要拿规格化之后的三个附加位再做一次判断。把对阶移位、尾数相加、前导零计数、规格化移位和舍入全部安排在一个周期内，关键路径会相当长；真实浮点加法器把它们分到不同流水级，这正是上一个专题说浮点运算“常常多周期”的具体原因。

\topic{专题：ALU 测试向量要覆盖边界值}
ALU 的测试最怕只使用随机正常值。随机测试当然有用，但边界值才能暴露进位、溢出、符号扩展、移位宽度和比较解释错误。一张合格的测试向量表至少应覆盖零、最大正数、最小负数、全 1、单 bit、交替 bit、进位跨字节、符号位翻转和移位距离边界。

\begin{longtable}{L{0.18\linewidth}L{0.26\linewidth}L{0.46\linewidth}}
\toprule
向量类型 & 示例 & 主要检查点 \\
\midrule
零值 & \code{0+0}，\code{0-0} & Zero 标志、无额外进位、写回不遗漏 \\\tablerowrule
无符号进位 & \code{0xFF+0x01} & 结果回绕和 CarryOut 同时正确 \\\tablerowrule
有符号溢出 & \code{0x7F+0x01} & Sign 为 1 不等于“合法负数”，Overflow 必须置位 \\\tablerowrule
最小负数 & \code{0x80-0x01}，\code{-128/-1} & 补码非对称范围和除法溢出 \\\tablerowrule
比较混用 & \code{0xFF < 0x01} & 无符号为假，有符号为真 \\\tablerowrule
移位边界 & \code{1<<0}，\code{1<<7}，\code{1<<8} & 移位距离截断或非法规则必须明确 \\\tablerowrule
多字进位 & \code{0x00FF+0x0001} & 低字节进位进入高字节 \\
\bottomrule
\end{longtable}

调试时建议把每个向量写成五列：输入、期望结果、期望标志、经过的 ALU 路径、提交时钟沿。这样才能区分“组合逻辑算错”“标志解释错”“控制线选错路径”和“写回时机错”。



本节练习要求把 bit 解释、硬件路径和控制信号联系起来思考。答案不应仅写算术结果，还要说明该结果来自哪条组合路径，在哪个时钟边界被保存，哪些控制线必须同时成立。











































































































































\begin{itemize}
\item 位级机制：补码取负可写成按位取反再加一；减法可复用加法器执行“B 取反、最低位进位为 1”。
\item 标志规则：Carry 适合无符号边界，Overflow 适合补码有符号边界；不能把两者混成同一个“溢出”。
\item ALU 层面：候选计算并行产生，选择器选出一个结果，标志位保存结果的特征。
\item 数据通路层面：为 \code{R3 = R1 + R2} 写读地址、ALU 输入选择、\code{alu_op}、写回选择、写地址、写使能。
\item 控制层面：每个周期哪些写使能为 1，哪些 mux 选择哪一路，ALU 做什么，都要能列成 trace。
\item 检查题：若 ALU 结果已经正确，但写回 mux 选择线接错，寄存器结果会怎样？参考答案：寄存器会保存被错误选择的来源，例如旧寄存器值、内存数据或外部输入，而不是 ALU 结果。若示波器、LED 或仿真已经证明 ALU 输出正确，排查重点应转向 \code{wb_sel}、写回 mux、写地址、\code{reg_write} 和控制器状态，而不是继续修改 ALU。
\end{itemize}

经典教材会把控制器和数据通路并排讲：数据通路说明“能走哪些路”，控制器说明“本周期允许哪条路”。两者合起来，才是可执行动作。
\begin{keyidea}
最小 CPU 不是“ALU 加寄存器”的粗线图，而是一组按周期提交的状态转换。每条指令都必须说明：本周期读哪些旧状态，组合路径经过哪些部件，哪些控制线有效，时钟沿会提交哪些新状态，异常或等待是否会阻止提交。
\end{keyidea}

\begin{pdfdiagram}
  \node[hw state, minimum width=1.7cm] (pc) at (0,0.75) {PC};
  \node[hw compute, minimum width=2.0cm] (imem) at (2.55,0.75) {指令存储};
  \node[hw state, minimum width=1.55cm] (ir) at (5.05,0.75) {IR};
  \node[hw control, minimum width=1.9cm] (ctrl) at (7.55,0.75) {控制器};
  \node[hw state, minimum width=2.0cm] (rf) at (2.55,-1.0) {寄存器堆};
  \node[hw compute, minimum width=1.7cm] (alu) at (5.05,-1.0) {ALU};
  \node[hw compute, minimum width=2.0cm] (wb) at (7.55,-1.0) {写回 mux};
  \draw[hw bus] (pc) -- node[above=2pt,hw label,fill=white,inner sep=1pt]{取指地址} (imem);
  \draw[hw bus] (imem) -- node[above=2pt,hw label,fill=white,inner sep=1pt]{编码 bit} (ir);
  \draw[hw bus] (ir) -- node[above=2pt,hw label,fill=white,inner sep=1pt]{字段} (ctrl);
  \draw[hw return] (ctrl.south) |- node[pos=0.25,right,hw label,fill=white,inner sep=1pt]{控制线} (5.05,-0.18) -| (alu.north);
  \draw[hw bus] (rf) -- node[above=2pt,hw label,fill=white,inner sep=1pt]{源操作数} (alu);
  \draw[hw bus] (alu) -- node[above=2pt,hw label,fill=white,inner sep=1pt]{结果/标志} (wb);
  \draw[hw return] (wb.east) -- ++(0.55,0) -- ++(0,-0.7) node[pos=0.55,right=2pt,hw label,fill=white,inner sep=1pt]{写回} -| (rf.south);
  \draw[hw return] (ctrl.north) |- (3.85,1.75) -| (pc.north) node[pos=0.2,above=2pt,hw label,fill=white,inner sep=1pt]{下一 PC};
  \node[hw label, text width=9.4cm] at (4.1,-2.3) {最小 CPU 的核心观察对象是 trace：PC 取指，IR 保存编码，控制器打开数据通路，寄存器和 PC 在时钟沿提交。};
\end{pdfdiagram}
\diagramnote{最小 CPU 总图。箭头表示一条指令从取指到写回的完整路径，反馈线表示 PC 和寄存器堆在时钟边界更新。}

\section{三、从动作编码到控制字组织}

若希望同一套数据通路在不同时间做不同动作，就需要把“这一次做什么”也表示成 bit。动作编码可以先写成最小形式：

\begin{lstlisting}
opcode | dst | src_a | src_b
\end{lstlisting}

\code{opcode} 表示动作类型，例如 ADD、SUB、AND；\code{dst} 表示写入哪个寄存器；\code{src\_a} 和 \code{src\_b} 表示读取哪两个寄存器。控制器读取这些字段后，把它们译成 ALU op、读地址、写地址、写使能和写回选择。编码本身仍然是硬件输入，不是文字命令。

若有一组动作编码连续放在存储结构中，硬件还需要知道当前取哪一条。这个位置由 PC 保存。PC 本质上也是寄存器，只是它保存的是“下一条动作编码的地址”。

\begin{longtable}{L{0.24\linewidth}L{0.52\linewidth}}
\toprule
部件 & 作用 \\
\midrule
PC & 保存下一条编码地址 \\\tablerowrule
指令存储器 & 根据 PC 输出当前编码 \\\tablerowrule
控制器 & 根据编码字段产生控制信号 \\\tablerowrule
寄存器堆 & 保存通用数据 \\\tablerowrule
ALU & 做整数和逻辑运算 \\\tablerowrule
选择器网络 & 选择 ALU 输入、写回来源和下一 PC \\
\bottomrule
\end{longtable}

顺序执行时，下一 PC 可以由加法器产生：

\begin{lstlisting}
pc_next = pc_current + instruction_size
\end{lstlisting}

如果每条编码占 4 个字节，\code{instruction\_size} 就是 4；如果是可变长度指令，下一 PC 的计算会更复杂。顺序执行时，PC 这个状态寄存器在时钟沿保存新的地址。

最小单周期模型可以这样走：

\begin{longtable}{L{0.18\linewidth}L{0.52\linewidth}L{0.18\linewidth}}
\toprule
阶段 & 硬件动作 & 边界 \\
\midrule
取指 & PC 送入指令存储器，得到编码 & 组合路径 \\\tablerowrule
译码 & 控制器产生 mux、ALU、写使能信号 & 组合路径 \\\tablerowrule
执行 & 寄存器堆读值，ALU 计算 & 组合路径 \\\tablerowrule
写回 & 结果在时钟沿写入寄存器或 PC & 状态边界 \\
\bottomrule
\end{longtable}

用一条普通运算编码走读：当前位置 PC=100，指令存储器在地址 100 处保存的是：

\begin{lstlisting}
ADD R1, R2   // R1 = R1 + R2
\end{lstlisting}

周期开始时，PC 输出稳定为 100。这个地址进入指令存储器，存储器经过地址译码后把这条编码输出。控制器看到 opcode=ADD，于是把 \code{alu\_op} 设为 ADD，把寄存器堆的两个读地址分别设为 R1 和 R2，把写地址设为 R1，把 \code{wb\_sel} 设为 ALU，把 \code{reg\_write} 设为 1，同时让下一 PC 选择 PC 加指令长度。

这些控制信号稳定后，寄存器堆的两个读端口输出 R1 和 R2 当前的值。它们经过 ALU 输入选择器进入 ALU，加法路径产生结果。结果再经过写回选择器，到达寄存器堆的写数据端口。另一边，PC 加法器准备下一条指令地址，PC 选择器把它送到 PC 的下一值输入。只有到时钟沿，两个状态变化才同时提交：R1 保存 ALU 结果，PC 保存下一条地址。

把这次走读填上具体数值，抽象的路径描述就变成了可逐项核对的数值关系。设 \code{R1=5}、\code{R2=7}、指令长 2 字节：

\begin{longtable}{L{0.16\linewidth}L{0.34\linewidth}L{0.42\linewidth}}
\toprule
时刻 & 路径上的值 & 核对要点 \\
\midrule
周期开始 & \code{PC=100} & 取指地址稳定 \\\tablerowrule
取指后 & \code{IR = ADD R1,R2} & 编码来自 \code{MEM[100]} \\\tablerowrule
译码后 & \code{alu\_op=ADD}，读地址 R1/R2，写地址 R1，\code{wb\_sel=ALU}，\code{reg\_write=1} & 控制线与 opcode 一致 \\\tablerowrule
读寄存器 & \code{read\_a=5}，\code{read\_b=7} & 寄存器堆当前值 \\\tablerowrule
执行后 & \code{ALU\_out=12}，写回 mux 输出 12 & 加法路径和写回选择正确 \\\tablerowrule
时钟沿 & \code{R1 <- 12}，\code{PC <- 102} & 两个状态同时提交，R2 不变 \\
\bottomrule
\end{longtable}

任何一格对不上，故障范围就锁死了：\code{read\_a} 不是 5，查寄存器堆读地址；\code{ALU\_out} 不是 12，查 \code{alu\_op} 和进位链；12 没到写数据口，查 \code{wb\_sel}；R1 没变或 PC 没变，查 \code{reg\_write}、\code{pc\_write} 和时钟沿本身。这就是 trace 比“检查一下”可靠的地方。

分支也可以归结为下一 PC 选择。顺序情况下，下一 PC 是 PC 加指令长度；条件分支则由 ALU 比较结果和选择器决定下一 PC 来源：

\begin{lstlisting}
branch_taken = condition_met AND branch_op
next_pc = mux(branch_taken, pc_sequential, branch_target)
\end{lstlisting}

分支的硬件本质是改变 PC 这个寄存器的下一值。下一周期取到哪条编码，完全由这个新 PC 决定。

\topic{把指令周期写成逐阶段的执行过程}经典教材讲 CPU 时，常把一条指令拆成取指、译码、执行、访存、写回和下一 PC 更新。拆分的目的是把每个阶段对应到真实硬件路径。最小单周期 CPU 可以在一个长周期内完成所有路径；多周期 CPU 可以把这些路径分到多个时钟；流水线 CPU 可以让多条指令在不同阶段重叠推进。三者的共同点是：每个阶段都要说明源状态、组合路径、目标状态和异常边界。

\begin{longtable}{L{0.14\linewidth}L{0.34\linewidth}L{0.28\linewidth}L{0.14\linewidth}}
\toprule
阶段 & 主要硬件动作 & 关键控制信号 & 提交边界 \\
\midrule
IF 取指 & PC 送到指令存储器或 I-cache，取回指令 bit & \code{pc_out}、\code{imem_read}、\code{ir_write} & IR 或取指缓冲 \\\tablerowrule
ID 译码 & opcode、寄存器号、立即数字段被解释成控制线 & \code{decode}、\code{imm_kind}、\code{read_addr} & 可组合，也可锁存 \\\tablerowrule
EX 执行 & ALU 做加减、逻辑、比较或地址形成 & \code{alu_op}、\code{a_sel}、\code{b_sel} & 输出和标志 \\\tablerowrule
MEM 访存 & 对数据存储器、cache 或 MMIO 发起读写事务 & \code{mem_read}、\code{mem_write}、\code{byte_en} & MDR 或外部响应 \\\tablerowrule
WB 写回 & ALU 结果或读回数据写入寄存器堆 & \code{wb_sel}、\code{reg_write} & 目标寄存器 \\\tablerowrule
PC 更新 & 选择顺序 PC、分支目标、跳转目标或异常入口 & \code{pc_sel}、\code{pc_write} & PC \\
\bottomrule
\end{longtable}

这张表也可以直接用作调试时的排查清单。若某条加载指令读到了错误数据，不应只问“内存为什么错”，而要逐项检查：PC 是否取到正确编码，译码是否识别为加载，ALU 是否形成正确地址，MEM 阶段是否访问正确目标，读回数据是否进入 MDR（memory data register，存储器数据寄存器——访存返回数据在进入写回路径之前的暂存，写回 mux 的“存储器数据”候选就是从它引出），WB 阶段是否选择了 MDR 而不是 ALUOut。这样的排查方式比凭感觉改控制器可靠得多。

\topic{控制字把“会执行指令”变成可接线信号}真正能搭出来的 CPU 必须有控制字。控制字是一组在某个周期稳定的控制线集合。它决定寄存器是否写入、ALU 选哪个操作、存储器是否读写、PC 从哪里取下一值。若控制字没有定义清楚，CPU 即使图上连通，也无法稳定执行程序。

\begin{longtable}{L{0.18\linewidth}L{0.2\linewidth}L{0.2\linewidth}L{0.2\linewidth}L{0.12\linewidth}}
\toprule
控制对象 & ADD & MOV（加载） & MOV（存储） & JZ \\
\midrule
ALU 操作 & 加法或指定算术 & 基址加偏移 & 基址加偏移 & 比较或测试零 \\\tablerowrule
寄存器读 & 两个源寄存器 & 基址寄存器 & 基址和数据寄存器 & 条件源寄存器 \\\tablerowrule
存储器读 & 无 & 有 & 无 & 无 \\\tablerowrule
存储器写 & 无 & 无 & 有 & 无 \\\tablerowrule
写回选择 & ALU 结果 & 存储器数据 & 无 & 无 \\\tablerowrule
PC 选择 & 顺序 PC & 顺序 PC & 顺序 PC & 条件目标或顺序 PC \\
\bottomrule
\end{longtable}

若把这些控制对象编码成 bit，可以得到类似下面的教学控制字。具体 bit 位不重要，重要的是每一位都能对应到实际连线或选择器输入。

\begin{lstlisting}
control word fields:
  pc_sel      // sequential, branch_target, jump_target, exception
  reg_write   // enable register-file write port
  wb_sel      // ALU, memory, pc_plus_size
  alu_op      // add, sub, and, or, compare
  alu_b_sel   // register or immediate
  mem_read    // request data read
  mem_write   // request data write
  byte_en     // selected byte lanes
  flags_write // update Zero/Carry/Overflow/Sign
\end{lstlisting}

把控制字应用到这台具体 ISA（编码见第二节）上，每条指令一行，全部 7 条覆盖：

{\setlength{\tabcolsep}{2pt}
\begin{longtable}{L{0.09\linewidth}L{0.1\linewidth}L{0.11\linewidth}L{0.11\linewidth}L{0.11\linewidth}L{0.07\linewidth}L{0.09\linewidth}L{0.09\linewidth}L{0.09\linewidth}}
\toprule
指令 & \code{alu\_op} & \code{a\_sel} & \code{b\_sel} & \code{wb\_sel} & \code{reg\_w} & \code{mem\_r} & \code{mem\_w} & \code{pc\_sel} \\
\midrule
\code{ADD} & ADD & 寄存器 & 寄存器 & ALU & 1 & 0 & 0 & 顺序 \\\tablerowrule
\code{MOV}（立即数） & 直送 & 立即数 & — & ALU & 1 & 0 & 0 & 顺序 \\\tablerowrule
\code{MOV}（加载） & ADD & 寄存器 & off4 扩展 & 存储器 & 1 & 1 & 0 & 顺序 \\\tablerowrule
\code{MOV}（存储） & ADD & 寄存器 & off4 扩展 & — & 0 & 0 & 1 & 顺序 \\\tablerowrule
\code{JZ} & 判零 & 寄存器 & — & — & 0 & 0 & 0 & Z?目标：顺序 \\\tablerowrule
\code{JMP} & — & — & — & — & 0 & 0 & 0 & 目标 \\\tablerowrule
\code{SHL} & 左移 & 寄存器 & imm4 & ALU & 1 & 0 & 0 & 顺序 \\
\bottomrule
\end{longtable}
}

这张表还能进一步写成布尔式。以 \code{reg\_write} 为例，它只在四条指令下为 1——\code{ADD}、立即数 \code{MOV}、加载 \code{MOV} 和 \code{SHL}；多条 \code{MOV} 在布尔式中用 opcode 值区分：

\begin{lstlisting}
reg_write = (op == ADD) OR (op == 0010) OR (op == 0011) OR (op == SHL)
mem_read  = (op == 0011)   // MOV load
mem_write = (op == 0100)   // MOV store
pc_sel    = (op == JMP) OR ((op == JZ) AND Zero) ? target : sequential
\end{lstlisting}

译码器到此就没有黑箱了：输入 opcode 的 4 个 bit，输出这张表里的每一格，且每一格都能在这张真值表和数据通路图之间互相指认。若某条指令行为异常，先在这张表上查控制字哪一格错了，再回数据通路图查那条线——排查顺序永远是先控制字、后路径。

\topic{硬连线控制与微程序控制是两种组织控制字的方法}控制器可以直接由 opcode、状态位和当前周期译码出控制线，这称为硬连线控制。它速度快，结构适合简单指令集和高频流水线，但修改指令行为需要改组合逻辑。另一种方法是把控制字存在控制存储器中，由微程序计数器逐条取出微指令；每条机器指令对应一段微操作序列。这种方法更像“硬件内部的小程序”，适合组织复杂指令、多周期总线访问和异常入口，但控制存储器和微序列会增加延迟。

\begin{longtable}{L{0.18\linewidth}L{0.34\linewidth}L{0.38\linewidth}}
\toprule
控制方式 & 适合场景 & 关键问题 \\
\midrule
硬连线控制 & 指令格式规整、周期少、目标频率高 & opcode、状态位和时序状态是否唯一决定控制线 \\\tablerowrule
微程序控制 & 指令复杂、多周期动作多、需要复用数据通路 & 微地址如何选择、分支微指令如何跳转、异常如何打断 \\\tablerowrule
混合方式 & 常见简单指令走快路径，复杂指令走微序列 & 快路径和微序列是否共享一致的提交边界 \\
\bottomrule
\end{longtable}

\begin{lstlisting}
// 注释（推进）：每个 u/T 行代表一个控制节拍，同一行内的控制信号并行生效。
// 注释（边界）：next 决定下一微地址；等待或错误时不得提前进入写回。
microinstruction sketch:
  alu_op | src_a | src_b | dst | mem_cmd | pc_cmd | next_micro_pc

load sequence:
  u0: IR  <- MEM[PC],       next = u1
  u1: A   <- R[base],       next = u2
  u2: ADR <- A + imm,       next = u3
  u3: MDR <- MEM[ADR],      next = u4 or wait
  u4: R[dst] <- MDR, PC <- PC+4, next = fetch
\end{lstlisting}

微程序写法有一个重要好处：它迫使作者说明每个周期保存什么中间状态。若某一步需要等待外部存储器，\code{next\_micro\_pc} 可以停留在当前微地址；若访问失败，微序列可以跳到异常入口。这样，等待、错误和普通写回不再混在一句“执行加载”里。

把这段微指令草图对应的硬件画出来，微程序控制器就是一台嵌在 CPU 内部的小取指机：

\begin{pdfdiagram}
  \node[hw state, minimum width=2.2cm, align=center] (upc) at (0,2.6) {微程序\\计数器 µPC};
  \node[hw memory, minimum width=2.0cm, align=center] (cmem) at (3.0,2.6) {控制\\存储器};
  \node[hw state, minimum width=2.0cm, align=center] (uir) at (6.0,2.6) {微指令\\寄存器};
  \node[hw compute, minimum width=2.6cm, align=center] (dp) at (6.0,0.4) {数据通路\\（控制线驱动）};
  \node[hw control, minimum width=2.2cm, align=center] (seq) at (1.8,0.4) {次微地址逻辑};
  \node[hw state, minimum width=1.4cm, align=center] (iro) at (4.3,0.4) {IR\\opcode};
  \draw[hw bus] (upc) -- node[hw label, above]{微地址} (cmem);
  \draw[hw bus] (cmem) -- node[hw label, above]{微指令} (uir);
  \draw[hw bus] (uir.south) -- node[hw label, right]{控制字段} (dp.north);
  \draw[hw return] (uir.south) -- (6.0,1.9) -- node[hw label, above, pos=0.55]{下一微地址字段} (1.8,1.9) -- (seq.north);
  \draw[hw return] (iro.west) -- node[hw label, above]{入口映射} (seq.east);
  \draw[hw return] (seq.west) -- node[hw label, below, pos=0.4]{次微地址} (-1.8,0.4) -- (-1.8,2.6) -- (upc.west);
\end{pdfdiagram}
\diagramnote{微程序控制器的结构骨架。µPC 保存当前微地址，控制存储器按它取出整条微指令，微指令寄存器把字段分成两路：控制字段直接驱动数据通路的各条控制线；下一微地址字段与 IR 的 opcode 入口映射一起进入次微地址逻辑，决定 µPC 的下一值。访存等待时令次微地址等于当前微地址原地停留，异常时改指异常微程序入口——控制规则全部写在控制存储器的表格内容里，数据通路本身不变。}

\topic{多周期 CPU 用微操作缩短最长组合路径}单周期 CPU 概念清楚，但周期必须长到容纳取指、译码、ALU、访存和写回的最坏组合路径。多周期 CPU 把同一条指令拆成若干微操作，让每个周期只完成一小段界限清楚的动作。例如加载指令可以拆为：PC 取指、译码读寄存器、ALU 形成地址、存储器读、写回寄存器。每一步之间用 IR、MDR、ALUOut 或临时寄存器保存中间结果。

在切分之前，先把这条决定周期长度的组合路径本身画出来：

\begin{pdfdiagram}
  \node[hw label] at (4.8,1.85) {最坏指令（加载）：五段组合路径在一拍内相继稳定};
  \node[hw compute, minimum width=1.6cm, align=center] (s1) at (0.8,0.7) {取指\\1 单位\\占 1/5};
  \node[hw compute, minimum width=1.6cm, align=center] (s2) at (2.8,0.7) {译码\\1 单位\\占 1/5};
  \node[hw compute, minimum width=1.6cm, align=center] (s3) at (4.8,0.7) {ALU\\1 单位\\占 1/5};
  \node[hw compute, minimum width=1.6cm, align=center] (s4) at (6.8,0.7) {访存\\1 单位\\占 1/5};
  \node[hw compute, minimum width=1.6cm, align=center] (s5) at (8.8,0.7) {写回\\1 单位\\占 1/5};
  \draw[hw bus] (s1) -- (s2);
  \draw[hw bus] (s2) -- (s3);
  \draw[hw bus] (s3) -- (s4);
  \draw[hw bus] (s4) -- (s5);
  % 周期长度标尺（刻度在各段右缘，对应累计延迟）
  \draw[line width=0.55pt] (0,-0.45) -- (9.6,-0.45);
  \draw[line width=0.55pt] (0,-0.45) -- (0,-0.3);
  \draw[line width=0.55pt] (1.6,-0.45) -- (1.6,-0.3);
  \draw[line width=0.55pt] (3.6,-0.45) -- (3.6,-0.3);
  \draw[line width=0.55pt] (5.6,-0.45) -- (5.6,-0.3);
  \draw[line width=0.55pt] (7.6,-0.45) -- (7.6,-0.3);
  \draw[line width=0.55pt] (9.6,-0.45) -- (9.6,-0.3);
  \node[hw label, below] at (0,-0.45) {0};
  \node[hw label, below] at (1.6,-0.45) {1};
  \node[hw label, below] at (3.6,-0.45) {2};
  \node[hw label, below] at (5.6,-0.45) {3};
  \node[hw label, below] at (7.6,-0.45) {4};
  \node[hw label, below] at (9.6,-0.45) {5};
  \node[hw label, anchor=west] at (9.75,-0.45) {时间单位};
  % 周期末提交沿
  \draw[hw return] (9.85,-0.2) -- (9.85,1.5);
  \node[hw label, anchor=east] at (9.7,1.4) {周期末时钟沿提交};
  \node[hw label] at (4.8,-1.1) {时钟周期必须 $\geq 5$ 个时间单位，由这条关键路径决定};
\end{pdfdiagram}
\diagramnote{单周期 CPU 的关键路径。取指、译码、ALU、访存、写回五段组合逻辑在一拍内串行稳定；按本章统一前提每段延迟 1 个时间单位、各占五分之一，最坏的加载指令要走完五段，故时钟周期必须取 5 个单位以上。多周期把这条链拆成 T0--T4 五拍，流水线把它隔成五级——两种组织都是在段边界处把这条长链切开。}

\begin{longtable}{L{0.12\linewidth}L{0.28\linewidth}L{0.3\linewidth}L{0.2\linewidth}}
\toprule
周期 & 微操作 & 打开的硬件路径 & 提交对象 \\
\midrule
T0 & \code{IR = MEM[PC]} & PC 到指令存储器，读响应到 IR & IR \\\tablerowrule
T1 & \code{A = R[rs], B = R[rt]} & 指令字段到寄存器堆读地址 & A/B 临时寄存器 \\\tablerowrule
T2 & \code{ALUOut = A + imm} & ALU 做地址形成或比较 & ALUOut/Flags \\\tablerowrule
T3 & \code{MDR = MEM[ALUOut]} & 地址到数据存储器，读响应到 MDR & MDR \\\tablerowrule
T4 & \code{R[rd] = MDR, PC = PC+4} & 写回 mux 和 PC 加法器 & 寄存器、PC \\
\bottomrule
\end{longtable}

这个 trace 说明，多周期设计把过长的硬件路径切成可计时、可观察、可插入等待的边界。若数据存储器慢，可以只让 T3 等待；若分支比较已经在 T2 得到结果，可以在 T2 或 T3 改写 PC；若发生异常，可以把 PC 选择器改到异常入口，而不是继续提交错误写回。

\topic{微操作记法把每条指令写成可核对的寄存器传输序列}前面的微指令草图和多周期 trace 已经用过 \code{IR = MEM[PC]} 这类写法，现在把这套记法定死，称为寄存器传输级记法。先交代一句：草图里的 \code{=} 与本节的 \code{<-} 含义相同，都表示时钟沿提交，下文统一写作 \code{<-}。核心约定只有三条：第一，箭头左侧是时钟沿要提交的状态元件，右侧读取的全部是本周期的旧值；第二，写在同一行里的多个传输在同一个时钟沿同时提交，行内没有先后；第三，行与行之间才是时间边界——多周期 CPU 里一行对应一拍，单周期 CPU 里整条指令写进一行，表示这些传输在同一拍内相继稳定、同一沿提交。于是 \code{R[d] <- R[d]+R[b]} 读作“时钟沿到来时，d 获得 d 与 b 旧值之和”，\code{MEM[R[a]+off] <- R[s]} 读作“时钟沿把 s 的旧值写入地址 R[a]+off 处”。

自然语言的“先取指、再执行”说不清哪些动作并发、哪些动作分拍，而硬件恰恰在这一点上不允许含糊。把指令写成传输序列之后，单周期与多周期之争就变成同一序列的两种排布——是一行写完，还是拆成五行、行间插入 IR、MDR、ALUOut 这些暂存元件——两种写法描述的是同一台机器。

\begin{longtable}{L{0.2\linewidth}L{0.4\linewidth}L{0.28\linewidth}}
\toprule
写法 & 含义 & 核对要点 \\
\midrule
\code{R[x]} & 寄存器堆中编号 x 的当前值 & 读到的是旧值 \\\tablerowrule
\code{MEM[a]} & 地址 a 处的存储字 & 地址须先形成 \\\tablerowrule
\code{A <- B} & 时钟沿把 B 提交给 A & 右侧全是旧值 \\\tablerowrule
同一行多个传输 & 同一时钟沿同时提交 & 行内无先后 \\\tablerowrule
\code{cond ? x : y} & 条件成立取 x，否则取 y & 条件须提前稳定 \\
\bottomrule
\end{longtable}

用这套记法，本章 ISA（编码见第二节）的四条代表指令各写成一行。注意本章 ROM 按字编址，顺序执行的下一 PC 是 \code{PC+1}：

\begin{lstlisting}
// 注释（推进）：每个 u/T 行代表一个控制节拍，同一行内的控制信号并行生效。
// 注释（边界）：next 决定下一微地址；等待或错误时不得提前进入写回。
ADD Rd, Rs:        IR <- ROM[PC];  R[d] <- R[d]+R[b];  PC <- PC+1
MOV Rd, [Ra+off4]: IR <- ROM[PC];  R[d] <- MEM[R[a]+sign_ext(off4)];  PC <- PC+1
MOV [Ra+off4], Rs: IR <- ROM[PC];  MEM[R[a]+sign_ext(off4)] <- R[s];  PC <- PC+1
JZ Rr, +off:       IR <- ROM[PC];  PC <- (R[r]==0) ? PC+1+sign_ext(off8) : PC+1
\end{lstlisting}

每一行都能翻成数值核对。以 \code{PC=2}、\code{R1=0xF000} 执行 \code{MOV R2, [R1+4]}（ROM 字 \code{0x3450}，编码见第二节）：本沿读旧值 \code{R[a]=0xF000}，地址形成 \code{0xF000+4=0xF004}，提交 \code{R2 <- MEM[0xF004]} 与 \code{PC <- 3}。以 \code{PC=3}、\code{R2=0} 执行 \code{JZ R2, +3}（\code{0x5403}）：条件成立，提交 \code{PC <- 3+1+3=7}；若 \code{R2=1}，同一行记法给出 \code{PC <- 4}。多周期写法只是把这些传输拆到 T0–T4 五行、行间用暂存元件保存中间值，与上一张多周期 trace 表逐行对应；单周期写法则把同一行内的传输理解为一条长组合路径上相继稳定的各段。

分支目标的形成路径值得单独画一张图，两路候选 PC 和选择条件可以逐项核对：

\begin{pdfdiagram}
  \node[hw state, minimum width=1.4cm] (pc) at (0,3.2) {PC};
  \node[hw compute, minimum width=1.5cm, align=center] (add1) at (2.2,3.2) {+1\\加法器};
  \node[hw state, minimum width=1.8cm, align=center] (irl) at (0,1.2) {IR 低位段\\off8 / off12};
  \node[hw compute, minimum width=1.7cm, align=center] (se) at (2.2,1.2) {符号扩展};
  \node[hw compute, minimum width=1.6cm, align=center] (addoff) at (4.6,1.2) {+off\\加法器};
  \node[hw control, minimum width=0.9cm] (zf) at (6.4,2.2) {Zero};
  \node[hw compute, minimum width=1.6cm, align=center] (pmux) at (8.0,2.2) {pc\\mux};
  \draw[hw bus] (pc) -- node[hw label, above]{本条地址} (add1);
  \draw[hw bus] (irl) -- (se);
  \draw[hw bus] (pc.south) -- (0,2.35) -- (4.6,2.35) -- (addoff.north);
  \draw[hw bus] (se) -- node[hw label, above]{sign\_ext} (addoff);
  \draw[hw bus] (add1.east) -- node[hw label, above, pos=0.6]{PC+1（顺序）} (8.0,3.2) -- (pmux.north);
  \draw[hw bus] (addoff.east) -- node[hw label, below, pos=0.65]{PC+1+off（目标）} (8.0,1.2) -- (pmux.south);
  \draw[hw return] (zf) -- (pmux);
  \draw[hw bus] (pmux.east) -- ++(0.7,0) |- (0,4.1) -- (pc.north);
  \node[hw label, above] at (4.5,4.1) {下一 PC};
\end{pdfdiagram}
\diagramnote{分支目标的地址形成。顺序值 PC+1 与目标值 PC+1+sign\_ext(off) 同时算好，由 Zero 选择一路作为下一 PC 提交；目标相对下一条指令计算，偏移按补码解释、可正可负。对照 ROM 程序：地址 3 的 \code{JZ R2, +3} 成立时目标是 $3+1+3=7$；地址 6 的 \code{JMP -5} 目标是 $6+1-5=2$；\code{JZ} 不成立时选择 PC+1，顺序取地址 4。}

\topic{同一程序在三种组织下的拍数对照}单周期、多周期和流水线是同一组硬件路径的三种时间安排。比较它们必须先固定前提：三种组织使用同一批部件——同一个指令存储、同一个寄存器堆、同一个 ALU、同一个数据存储；设每段组合路径（取指、译码、执行、访存、写回）延迟相同，记为 1 个时间单位，并忽略流水寄存器自身的建立延迟。在此前提下，单周期机的时钟周期必须容纳最坏指令（加载：五段串行），即 5 个单位；多周期和流水线的周期由最长一段决定，即 1 个单位。

考察下面这个 6 条指令的小程序（取自“极小 ISA 与可接线数据通路”一节 ROM 程序的前 6 条，设按键已按下，GPIO\_IN 读出 1，因此 \code{JZ} 条件不成立、顺序执行）：

\begin{lstlisting}
0: MOV R1, 0xF0     // R1 = 0x00F0
1: SHL R1, 8        // R1 = 0xF000 (MMIO base)
2: MOV R2, [R1+4]   // R2 = MEM[0xF004], GPIO_IN reads 1
3: JZ R2, +3        // R2 != 0, not taken
4: MOV R3, 1        // R3 = 1
5: MOV [R1+0], R3   // MEM[0xF000] = 1, LED on
\end{lstlisting}

单周期组织最直接：每条 1 拍，共 $6\times1=6$ 拍，每拍 5 个时间单位，总时间 $6\times5=30$。

多周期组织按指令分流。公共的取指、译码各占 1 拍；立即数 \code{MOV}、\code{SHL} 与 \code{JZ} 这类非访存指令再走执行、写回，各 4 拍；加载与存储多一拍访存，各 5 拍。总拍数 $4+4+5+4+4+5=26$，总时间 26 个单位，平均每条约 4.3 拍。

理想 5 级流水（取指、译码、执行、访存、写回各一级）先看无冒险情形：$5+6-1=10$ 拍——前 4 拍是填充，此后稳态每拍完成 1 条。再逐对检查相关：\code{SHL} 读 \code{R1} 紧跟立即数 \code{MOV}，加载的基址又读 \code{R1}，这两对靠 ALU 出口到入口的转发解决，不停顿；\code{JZ} 读 \code{R2}，而 \code{R2} 来自上一条加载，数据要到访存拍末才就绪，这是 load-use 冒险，即使有转发也要停 1 拍；\code{JZ} 不跳且按“猜不跳”预取，猜中无冲刷——若 \code{R2=0} 真跳，已在取指、译码级的两条预取作废，要再加 2 拍。本场景总拍数 $10+1=11$，总时间 11 个单位。

\begin{longtable}{L{0.14\linewidth}L{0.14\linewidth}L{0.18\linewidth}L{0.4\linewidth}}
\toprule
组织 & 总拍数 & 折合时间 & 关键问题 \\
\midrule
单周期 & 6 拍 & $6\times5=30$ & 周期被最慢的加载路径拖长 \\\tablerowrule
多周期 & 26 拍 & $26\times1=26$ & 周期短，但每条仍要排满 4–5 拍，平均每条约 4.3 拍 \\\tablerowrule
理想流水 & 11 拍 & $11\times1=11$ & 填充 4 拍、load-use 停 1 拍；程序越长平均拍数越接近 1 \\
\bottomrule
\end{longtable}

把三种组织画到同一条时间轴上，拍数与周期长度的差别直接可见：

\begin{pdfdiagram}
  % ===== 单周期：6 拍，每拍 5 个时间单位 =====
  \node[hw label] at (4.5,9.55) {每拍 5 个时间单位（周期按最慢的加载路径定）};
  \node[hw label, anchor=east] at (-0.2,8.8) {单周期};
  \draw[BookTeal, line width=0.9pt] (0,8.4) rectangle (1.5,9.2);
  \node[font=\tiny\sffamily] at (0.75,8.8) {I1};
  \draw[BookTeal, line width=0.9pt] (1.5,8.4) rectangle (3.0,9.2);
  \node[font=\tiny\sffamily] at (2.25,8.8) {I2};
  \draw[BookTeal, line width=0.9pt] (3.0,8.4) rectangle (4.5,9.2);
  \node[font=\tiny\sffamily] at (3.75,8.8) {I3};
  \draw[BookTeal, line width=0.9pt] (4.5,8.4) rectangle (6.0,9.2);
  \node[font=\tiny\sffamily] at (5.25,8.8) {I4};
  \draw[BookTeal, line width=0.9pt] (6.0,8.4) rectangle (7.5,9.2);
  \node[font=\tiny\sffamily] at (6.75,8.8) {I5};
  \draw[BookTeal, line width=0.9pt] (7.5,8.4) rectangle (9.0,9.2);
  \node[font=\tiny\sffamily] at (8.25,8.8) {I6};
  \node[hw label, anchor=west] at (9.15,8.8) {总时间 30};
  % ===== 多周期：4/5 拍分流，共 26 拍 =====
  \node[hw label, anchor=east] at (-0.2,7.0) {多周期};
  \draw[BookTeal, line width=0.9pt] (0,6.6) rectangle (1.2,7.4);
  \node[font=\tiny\sffamily, align=center] at (0.6,7.0) {I1\\4 拍};
  \draw[BookTeal, line width=0.9pt] (1.2,6.6) rectangle (2.4,7.4);
  \node[font=\tiny\sffamily, align=center] at (1.8,7.0) {I2\\4 拍};
  \draw[BookTeal, line width=0.9pt] (2.4,6.6) rectangle (3.9,7.4);
  \node[font=\tiny\sffamily, align=center] at (3.15,7.0) {I3 加载\\5 拍};
  \draw[BookTeal, line width=0.9pt] (3.9,6.6) rectangle (5.1,7.4);
  \node[font=\tiny\sffamily, align=center] at (4.5,7.0) {I4\\4 拍};
  \draw[BookTeal, line width=0.9pt] (5.1,6.6) rectangle (6.3,7.4);
  \node[font=\tiny\sffamily, align=center] at (5.7,7.0) {I5\\4 拍};
  \draw[BookTeal, line width=0.9pt] (6.3,6.6) rectangle (7.8,7.4);
  \node[font=\tiny\sffamily, align=center] at (7.05,7.0) {I6 存储\\5 拍};
  \node[hw label, anchor=west] at (7.95,7.0) {总时间 26};
  % ===== 五级流水：11 拍，含 1 拍 load-use 气泡 =====
  \node[hw label] at (1.9,5.85) {F=取指 D=译码 E=执行 M=访存 W=写回};
  \node[hw label, rotate=90] at (-1.35,4.14) {五级流水};
  \node[hw label, anchor=east] at (-0.2,5.29) {I1};
  \node[hw label, anchor=east] at (-0.2,4.83) {I2};
  \node[hw label, anchor=east] at (-0.2,4.37) {I3};
  \node[hw label, anchor=east] at (-0.2,3.91) {I4};
  \node[hw label, anchor=east] at (-0.2,3.45) {I5};
  \node[hw label, anchor=east] at (-0.2,2.99) {I6};
  % I1
  \draw[BookTeal, line width=0.7pt] (0,5.08) rectangle (0.3,5.5);
  \node[font=\tiny\sffamily] at (0.15,5.29) {F};
  \draw[BookTeal, line width=0.7pt] (0.3,5.08) rectangle (0.6,5.5);
  \node[font=\tiny\sffamily] at (0.45,5.29) {D};
  \draw[BookTeal, line width=0.7pt] (0.6,5.08) rectangle (0.9,5.5);
  \node[font=\tiny\sffamily] at (0.75,5.29) {E};
  \draw[BookTeal, line width=0.7pt] (0.9,5.08) rectangle (1.2,5.5);
  \node[font=\tiny\sffamily] at (1.05,5.29) {M};
  \draw[BookTeal, line width=0.7pt] (1.2,5.08) rectangle (1.5,5.5);
  \node[font=\tiny\sffamily] at (1.35,5.29) {W};
  % I2
  \draw[BookTeal, line width=0.7pt] (0.3,4.62) rectangle (0.6,5.04);
  \node[font=\tiny\sffamily] at (0.45,4.83) {F};
  \draw[BookTeal, line width=0.7pt] (0.6,4.62) rectangle (0.9,5.04);
  \node[font=\tiny\sffamily] at (0.75,4.83) {D};
  \draw[BookTeal, line width=0.7pt] (0.9,4.62) rectangle (1.2,5.04);
  \node[font=\tiny\sffamily] at (1.05,4.83) {E};
  \draw[BookTeal, line width=0.7pt] (1.2,4.62) rectangle (1.5,5.04);
  \node[font=\tiny\sffamily] at (1.35,4.83) {M};
  \draw[BookTeal, line width=0.7pt] (1.5,4.62) rectangle (1.8,5.04);
  \node[font=\tiny\sffamily] at (1.65,4.83) {W};
  % I3（加载）
  \draw[BookTeal, line width=0.7pt] (0.6,4.16) rectangle (0.9,4.58);
  \node[font=\tiny\sffamily] at (0.75,4.37) {F};
  \draw[BookTeal, line width=0.7pt] (0.9,4.16) rectangle (1.2,4.58);
  \node[font=\tiny\sffamily] at (1.05,4.37) {D};
  \draw[BookTeal, line width=0.7pt] (1.2,4.16) rectangle (1.5,4.58);
  \node[font=\tiny\sffamily] at (1.35,4.37) {E};
  \draw[BookTeal, line width=0.7pt] (1.5,4.16) rectangle (1.8,4.58);
  \node[font=\tiny\sffamily] at (1.65,4.37) {M};
  \draw[BookTeal, line width=0.7pt] (1.8,4.16) rectangle (2.1,4.58);
  \node[font=\tiny\sffamily] at (1.95,4.37) {W};
  % I4（JZ）：气泡 × 在第 6 拍
  \draw[BookTeal, line width=0.7pt] (0.9,3.7) rectangle (1.2,4.12);
  \node[font=\tiny\sffamily] at (1.05,3.91) {F};
  \draw[BookTeal, line width=0.7pt] (1.2,3.7) rectangle (1.5,4.12);
  \node[font=\tiny\sffamily] at (1.35,3.91) {D};
  \draw[BookAccent, line width=0.9pt] (1.5,3.7) rectangle (1.8,4.12);
  \node[font=\tiny\sffamily] at (1.65,3.91) {$\times$};
  \draw[BookTeal, line width=0.7pt] (1.8,3.7) rectangle (2.1,4.12);
  \node[font=\tiny\sffamily] at (1.95,3.91) {E};
  \draw[BookTeal, line width=0.7pt] (2.1,3.7) rectangle (2.4,4.12);
  \node[font=\tiny\sffamily] at (2.25,3.91) {M};
  \draw[BookTeal, line width=0.7pt] (2.4,3.7) rectangle (2.7,4.12);
  \node[font=\tiny\sffamily] at (2.55,3.91) {W};
  % I5：第 6 拍滞留取指（空心格）
  \draw[BookTeal, line width=0.7pt] (1.2,3.24) rectangle (1.5,3.66);
  \node[font=\tiny\sffamily] at (1.35,3.45) {F};
  \draw[BookRule, line width=0.5pt] (1.5,3.24) rectangle (1.8,3.66);
  \draw[BookTeal, line width=0.7pt] (1.8,3.24) rectangle (2.1,3.66);
  \node[font=\tiny\sffamily] at (1.95,3.45) {D};
  \draw[BookTeal, line width=0.7pt] (2.1,3.24) rectangle (2.4,3.66);
  \node[font=\tiny\sffamily] at (2.25,3.45) {E};
  \draw[BookTeal, line width=0.7pt] (2.4,3.24) rectangle (2.7,3.66);
  \node[font=\tiny\sffamily] at (2.55,3.45) {M};
  \draw[BookTeal, line width=0.7pt] (2.7,3.24) rectangle (3.0,3.66);
  \node[font=\tiny\sffamily] at (2.85,3.45) {W};
  % I6
  \draw[BookTeal, line width=0.7pt] (1.8,2.78) rectangle (2.1,3.2);
  \node[font=\tiny\sffamily] at (1.95,2.99) {F};
  \draw[BookTeal, line width=0.7pt] (2.1,2.78) rectangle (2.4,3.2);
  \node[font=\tiny\sffamily] at (2.25,2.99) {D};
  \draw[BookTeal, line width=0.7pt] (2.4,2.78) rectangle (2.7,3.2);
  \node[font=\tiny\sffamily] at (2.55,2.99) {E};
  \draw[BookTeal, line width=0.7pt] (2.7,2.78) rectangle (3.0,3.2);
  \node[font=\tiny\sffamily] at (2.85,2.99) {M};
  \draw[BookTeal, line width=0.7pt] (3.0,2.78) rectangle (3.3,3.2);
  \node[font=\tiny\sffamily] at (3.15,2.99) {W};
  \node[hw label, anchor=west] at (3.5,4.14) {总时间 11};
  \node[hw label] at (1.9,2.52) {$\times$ = load-use 气泡（停 1 拍）};
  % ===== 共用时间轴 =====
  \draw[line width=0.55pt] (0,2.2) -- (9.0,2.2);
  \draw[line width=0.55pt] (0,2.2) -- (0,2.35);
  \draw[line width=0.55pt] (1.5,2.2) -- (1.5,2.35);
  \draw[line width=0.55pt] (3.0,2.2) -- (3.0,2.35);
  \draw[line width=0.55pt] (4.5,2.2) -- (4.5,2.35);
  \draw[line width=0.55pt] (6.0,2.2) -- (6.0,2.35);
  \draw[line width=0.55pt] (7.5,2.2) -- (7.5,2.35);
  \draw[line width=0.55pt] (9.0,2.2) -- (9.0,2.35);
  \node[hw label, below] at (0,2.2) {0};
  \node[hw label, below] at (1.5,2.2) {5};
  \node[hw label, below] at (3.0,2.2) {10};
  \node[hw label, below] at (4.5,2.2) {15};
  \node[hw label, below] at (6.0,2.2) {20};
  \node[hw label, below] at (7.5,2.2) {25};
  \node[hw label, below] at (9.0,2.2) {30};
  \node[hw label, anchor=west] at (9.15,2.2) {时间单位};
\end{pdfdiagram}
\diagramnote{同一程序在三种组织下的时间轴（横轴共用同一时间单位）。单周期每条 1 拍、每拍 5 个单位，共 $6\times5=30$；多周期非访存指令 4 拍、加载与存储 5 拍，共 $4+4+5+4+4+5=26$；五级流水理想情形 10 拍，\code{JZ} 读 \code{R2} 撞上上一条加载的结果，插入 1 拍 load-use 气泡，共 11。流水行里 I4 的 $\times$ 就是这个气泡：I3 的数据要到访存拍末才就绪，I4 的执行拍只能推迟一格，I5 第 6 拍的空心格是随之滞留的取指。}

三个数字 30、26、11 说明吞吐率来自哪里。单周期把整条路径拉长成一拍，多周期把拍切短但每条指令仍要排满几拍，流水线则让多条指令同时占据不同段。真实机器里各段延迟并不相等、一次访存可能等几十拍、流水寄存器本身也有延迟，但只要比较的前提仍是同一组硬件路径，“切短周期、重叠执行”这个方向性结论就稳定成立。

\section{四、控制的基本元件：寄存器接口与状态机}

\topic{任何设备 = 一组寄存器 + 一台状态机}

把网卡、磁盘控制器、UART、定时器、GPU 命令处理器切开，结构惊人地一致：面向 CPU 的是一组寄存器，寄存器后面是一台状态机，状态机驱动设备的引脚或内部数据通路。寄存器按角色分三类：

- \hwkey{控制寄存器}（Control Register, CR）：软件写入，作为状态机的输入。写值告诉状态机做什么——启动、停止、选模式、开中断。
- \hwkey{状态寄存器}（Status Register, SR）：硬件写入、软件读出，是状态机对外的输出。读值告诉软件状态机现在在哪——忙不忙、有没有错、数据到了没有。
- \hwkey{数据寄存器}（Data Register, DR）：数据通路的出入口。写它把数据推入发送 FIFO，读它从接收 FIFO 弹出数据。数据寄存器通常有副作用：读一次少一个元素，写一次多一个元素。

这组寄存器怎样被物理地址选中？答案与最小 CPU 章的指令译码器是同一种电路：那里是 opcode 进译码器、控制线出；这里是物理地址高位进译码器、片选出。地址高位与设备被分配的地址区间比较，命中则产生片选使能该设备的寄存器堆，地址低位直接作为寄存器号在堆内选择——经典芯片与平台案例提到的 74LS138 三八译码器就是这套机制在分立元件时代的形态，今天它集成在根复合体和每个设备的地址译码逻辑里。软件一侧看到的对应物是驱动的 \code{ioremap}/\code{readl}/\code{writel}：\code{ioremap} 把这段物理区间映射进内核虚拟地址，并把它建立为不可缓存（UC）的映射（属性由 PAT/MTRR 保证）；\code{readl}/\code{writel} 则提供 volatile 语义与必要的顺序、屏障保证，使访问真正发到设备。

\begin{pdfdiagram}[x=1cm,y=1cm]
  % 地址译码器
  \node[hw compute, minimum width=2.4cm, minimum height=1.0cm] (dec) at (1.7,2.8) {地址译码器\\（区间比较）};
  \draw[hw bus] (-0.7,2.8) -- (0.5,2.8);
  \node[hw label, font=\tiny, anchor=south] at (-0.35,2.86) {地址高位};
  % 寄存器堆
  \node[hw state, minimum width=3.2cm, minimum height=2.6cm] (rf) at (5.0,0.6) {};
  \draw[BookTeal, line width=0.4pt] (3.7,1.25) rectangle (6.3,1.75);
  \node[hw label, font=\tiny] at (5.0,1.5) {控制寄存器 CR};
  \draw[BookTeal, line width=0.4pt] (3.7,0.45) rectangle (6.3,0.95);
  \node[hw label, font=\tiny] at (5.0,0.7) {状态寄存器 SR};
  \draw[BookTeal, line width=0.4pt] (3.7,-0.35) rectangle (6.3,0.15);
  \node[hw label, font=\tiny] at (5.0,-0.1) {数据寄存器 / FIFO};
  % 片选
  \draw[hw bus] (2.9,2.8) -- (3.8,2.8) -- (3.8,1.9);
  \node[hw label, font=\tiny, anchor=south] at (3.35,2.9) {片选};
  % 低位地址 / 数据
  \draw[hw bus] (-0.7,1.2) -- (3.4,1.2);
  \node[hw label, font=\tiny, anchor=south] at (0.5,1.28) {地址低位（寄存器号）};
  \draw[hw bus] (-0.7,0.5) -- (3.4,0.5);
  \node[hw label, font=\tiny, anchor=south] at (0.5,0.58) {写数据};
  \draw[hw bus] (3.4,-0.1) -- (-0.7,-0.1);
  \node[hw label, font=\tiny, anchor=north] at (0.5,-0.18) {读数据};
  % 状态机
  \node[hw control, minimum width=2.6cm, minimum height=2.6cm] (fsm) at (9.3,0.6) {设备状态机};
  \draw[hw bus] (6.6,1.5) -- (8.0,1.5);
  \node[hw label, font=\tiny, anchor=south] at (7.3,1.56) {控制位};
  \draw[hw return] (8.0,0.7) -- (6.6,0.7);
  \node[hw label, font=\tiny, anchor=south] at (7.3,0.76) {状态位};
  \draw[hw bus] (6.6,-0.1) -- (8.0,-0.1);
  \node[hw label, font=\tiny, anchor=south] at (7.3,-0.04) {发送数据};
  \draw[hw bus] (8.0,-0.45) -- (6.6,-0.45);
  \node[hw label, font=\tiny, anchor=north] at (7.3,-0.51) {接收数据};
  % 引脚
  \draw[hw bus] (10.6,0.9) -- (11.4,0.9);
  \node[hw label, font=\tiny, anchor=south] at (11.0,0.96) {输出引脚};
  \draw[hw bus] (11.4,0.25) -- (10.6,0.25);
  \node[hw label, font=\tiny, anchor=north] at (11.0,0.19) {输入引脚};
  % IRQ
  \draw[hw return] (9.3,1.9) -- (9.3,2.9);
  \node[hw label, font=\tiny, anchor=west] at (9.45,2.4) {IRQ};
\end{pdfdiagram}
\diagramnote{设备寄存器接口。地址高位经译码器产生片选选中整台设备，地址低位在寄存器堆内选择寄存器，与最小 CPU 章的 opcode 译码是同一种电路。CR 是状态机的输入，SR 是状态机的输出，数据寄存器/FIFO 连接状态机内部的数据通路，IRQ 是状态机向 CPU 方向的唯一主动信号。}

\topic{寄存器接口的读写约定：为什么会有这些奇怪规则}

翻开任何一份设备数据手册，寄存器位描述里总有一列读写属性，其中不少看起来不合直觉：有的位“读了就清零”，有的位“写 1 清零、写 0 无效”，有的位“置位后由硬件自己清回 0”。这些约定不是设计者的癖好，全部来自同一个结构性事实：\hwkey{寄存器有两个并发写者}——软件和硬件。普通内存只有一个写者（程序），设备寄存器却是软件与设备状态机共享的变量，竞态随时可能发生。

\begin{longtable}{L{0.21\linewidth}L{0.36\linewidth}L{0.36\linewidth}}
\toprule
约定 & 行为 & 为什么存在 \\
\midrule
read-to-clear & 读该寄存器后某些位自动清零 & 中断事件位：读走事件即表示“已取走”，硬件随即清位准备记录下一事件；免去了软件再写一次清除 \\\tablerowrule
W1C（写 1 清零） & write-1-to-clear：写 1 的位清零，写 0 的位不变 & 让软件只清自己处理的位，不清并发到达的新事件位 \\\tablerowrule
自清零（self-clearing） & 软件写 1 启动动作，硬件在动作开始后自动清回 0 & 启动位表示“请求”而非“状态”：软件读回 0 即知硬件已接手，避免软件忘清导致重复启动 \\\tablerowrule
只读（RO）/ 只写（WO） & 单向寄存器 & 状态位被软件写会制造假象；命令位被软件读回没有物理来源——索性在硬件上禁止 \\
\bottomrule
\end{longtable}

W1C 值得展开，因为它直接回应竞态。典型错误是“读-改-写”：软件读出 \code{0x05}（位 0 和位 2 有事件），处理完位 0 后清掉该位得到 \code{0x04}，再写回——若在读与写之间硬件刚置起了位 1，这次写回会把新事件位又抹回 0，事件被无声丢失。W1C 从根上消除这个窗口：写 \code{0x01} 只清位 0，写 0 的位硬件不理会，于是软件的清除动作与硬件的置位动作在位一级互不干涉，不需要任何锁。

\begin{warningbox}
把设备状态寄存器当普通变量做“读-改-写”是驱动开发中反复出现的缺陷来源。C 里的 \code{*p |= FLAG} 编译出来正是读-改-写三条指令，对 W1C 寄存器执行 \code{|=} 会把读到的所有事件位全部清掉（每个读到的 1 都被写回）。正确做法是直接把要清的位作为常量写入：\code{*p = FLAG}。另外，\code{volatile} 只阻止编译器优化掉这次访问，它管不了硬件并发——约定是硬件给的，遵守靠软件。
\end{warningbox}

\topic{轮询还是中断：一个代价模型}

设备完成一件事之后，软件有两种发现方式：反复读状态寄存器（轮询），或等设备发中断。两者各有确定的代价，选择是一个定量问题。

轮询的代价是持续的 CPU 占用。每次查询是一次 MMIO 读：读不能走 cache，必须发到设备再返回，PCIe 设备上一次 MMIO 读的往返在几百纳秒到 1 微秒量级。若每 $P$ 秒查询一次，CPU 在查询上的占用率约为 $T_q/P$（$T_q$ 是单次查询耗时），而事件被发现的最大延迟是 $P$。$P$ 取得小，发现得快但占用率线性上升；$P$ 取得大，省 CPU 但响应慢。中断的代价是每次事件一笔固定开销：CPU 保存现场、经向量表分发、进入服务程序，再恢复现场返回，全程在微秒量级，事件稀疏时几乎为零成本，但事件密集到接近这笔开销的倒数时，CPU 会被中断进入/退出本身耗尽——这正是网络设备在中断风暴下吞吐反而下降的原因。

\begin{longtable}{L{0.16\linewidth}L{0.38\linewidth}L{0.38\linewidth}}
\toprule
维度 & 轮询 & 中断 \\
\midrule
CPU 开销 & 与事件无关，$T_q/P$ 持续占用 & 与事件率成正比，每次约微秒量级 \\\tablerowrule
发现延迟 & 最坏一个查询周期 $P$，可预期 & 进入延迟加调度抖动，波动较大 \\\tablerowrule
密集事件 & 开销不随事件增长，反而划算 & 每次事件一笔固定开销，可能耗尽 CPU \\\tablerowrule
稀疏事件 & 绝大多数查询空手而归 & 几乎零成本 \\
\bottomrule
\end{longtable}

工程上的选择规则因此是：事件间隔远小于中断开销时轮询（高性能存储、用户态网络栈的 busy-poll），事件稀疏且到达时间不可预知时中断（键盘、串口），突发型负载用混合策略——Linux 网卡的 NAPI 是教科书式的例子：第一个包触发中断，服务程序关掉中断改为轮询，把积压的包一批收完，队列空了再重新开中断。中断负责“叫醒”，轮询负责“趁醒着多干”，两种机制的切换点由队列是否腾空这个状态决定——这本身就是一个小小的反馈环。

\topic{两个经典例子：8254 PIT 与 16550 UART}

\textbf{8254 可编程间隔定时器}（Programmable Interval Timer, PIT）把前面所有概念集于一身。它内含三个独立的 16 位递减计数器，共享一个 \hwevidence{1.193182 MHz} 的输入时钟；每个计数器装入初值 $N$ 后逐拍递减，数到 0 时按所选方式动作，并在方式 2/3 下自动重装初值。方式 3（方波发生器）下输出引脚 OUT 产生周期为 $N$ 个输入时钟的方波——分频器就是这么来的。IBM PC 把计数器 0 的输出接到 8259 的 IR0，成为操作系统的时钟节拍来源：初值 0 按 65536 计，$1193182/65536 \approx 18.2$ Hz 的节拍是 PC 兼容机沿用多年的传统。

\begin{longtable}{L{0.14\linewidth}L{0.24\linewidth}L{0.55\linewidth}}
\toprule
端口 & 寄存器 & 作用 \\
\midrule
\code{0x40} & 计数器 0 & 读：当前计数值；写：初值低/高字节 \\\tablerowrule
\code{0x41} & 计数器 1 & 同上（PC 上曾用于 DRAM 刷新节拍） \\\tablerowrule
\code{0x42} & 计数器 2 & 同上（PC 上接扬声器） \\\tablerowrule
\code{0x43} & 控制字寄存器（WO） & 选择计数器、读写顺序、工作方式、数制 \\
\bottomrule
\end{longtable}

控制字（写 \code{0x43}）的高两位选计数器，中间两位选读写顺序（先低字节后高字节；00 为锁存命令，即下文的计数值快照），接下来三位选方式 0--5，最低位选二进制或 BCD（Binary-Coded Decimal，二-十进制编码）计数。把计数器 0 配成方式 3、输出约 1 kHz 方波的完整序列只有六条指令（x86 端口 I/O，\code{out dx, al} 中端口地址经 \code{dx} 给出）：

\begin{lstlisting}
mov al, 0x36        ; 计数器0, 先低后高, 方式3, 二进制
out 0x43, al
mov al, 0xA9        ; 初值低字节
out 0x40, al
mov al, 0x04        ; 初值高字节: 0x04A9 = 1193
out 0x40, al        ; 1193182 / 1193 ≈ 1000 Hz
\end{lstlisting}

8254 还有一个细节恰好印证 1.2 节的并发问题：计数器在不停地减，软件读 16 位计数值要分两次读，两次读之间低字节可能向高字节借位，拼出来的值就是错的。解法是锁存命令——先向 \code{0x43} 发一条锁存命令，硬件把当前计数值瞬时复制到一个保持寄存器，软件再分两字节读这份快照。硬件更新与软件读取的竞态，由硬件提供一个原子快照来化解。

\begin{classicnote}
1.193182 MHz 这个古怪的频率来自 IBM PC 的 14.31818 MHz 晶体——它是 NTSC 彩色副载波 3.579545 MHz 的 4 倍，同一颗晶体既供显示又供 CPU 外设时钟；8254 的输入是它除以 12。经典芯片与平台案例讲过这种“一颗晶体供全机”的时代，HPET 出现之前的二十多年里，操作系统的时间观念就建立在这 18.2 Hz 的节拍上。
\end{classicnote}

\textbf{16550 UART} 是串口控制器，把 CPU 的并行字节与串行线上的位流互相转换（经典芯片与平台案例已在平台清单里提到它，这里拆寄存器）。它以 8 个端口偏移呈现一组寄存器，多数偏移上读和写是两个不同的寄存器：

\begin{longtable}{L{0.10\linewidth}L{0.10\linewidth}L{0.30\linewidth}L{0.36\linewidth}}
\toprule
偏移 & DLAB & 读 & 写 \\
\midrule
0 & 0 & RBR：接收缓冲（RO，读了就取走） & THR：发送保持（WO） \\\tablerowrule
0 & 1 & DLL：除数低字节 & DLL：除数低字节 \\\tablerowrule
1 & 0 & IER：中断使能 & IER：中断使能 \\\tablerowrule
1 & 1 & DLH：除数高字节 & DLH：除数高字节 \\\tablerowrule
2 & -- & IIR：中断标识（RO，读即应答） & FCR：FIFO 控制（WO） \\\tablerowrule
3 & -- & LCR：线路控制（位 7 即 DLAB） & LCR：线路控制 \\\tablerowrule
4 & -- & MCR：调制解调器控制 & MCR：调制解调器控制 \\\tablerowrule
5 & -- & LSR：线路状态（RO） & -- \\\tablerowrule
6 & -- & MSR：调制解调器状态（RO） & -- \\\tablerowrule
7 & -- & SCR：暂存 & SCR：暂存 \\
\bottomrule
\end{longtable}

这张表本身就是 1.2 节约定的陈列柜：RBR 是 read-to-clear 的数据寄存器，IIR 是 read-to-acknowledge 的状态寄存器，THR 只写、LSR 只读。DLAB（Divisor Latch Access Bit，除数锁存访问位，LCR 的位 7）是另一种约定——偏移 0/1 各有两个目标寄存器，地址位不够用了，就拿一个模式位当“第 9 位地址”：DLAB=1 时访问除数锁存器，DLAB=0 时访问数据与中断使能。波特率由除数决定：输入时钟 1.8432 MHz，先固定除以 16 再除以除数锁存器的值，除数为 1 即 \hwevidence{115200 baud}。初始化 115200 8N1 并打开 FIFO 的序列：

\begin{lstlisting}
; COM1 = 0x3F8
mov dx, 0x3FB       ; LCR
mov al, 0x80        ; DLAB=1
out dx, al
mov dx, 0x3F8       ; DLL
mov al, 1           ; 除数 1 -> 115200 baud
out dx, al
mov dx, 0x3F9       ; DLH = 0
xor al, al
out dx, al
mov dx, 0x3FB       ; LCR: DLAB=0, 8 数据位/无校验/1 停止位
mov al, 0x03
out dx, al
mov dx, 0x3FA       ; FCR: 使能 FIFO, 清收发 FIFO, 触发水位 14
mov al, 0xC7
out dx, al
\end{lstlisting}

发送时的轮询循环直接对应 1.3 节的代价模型——每次 \code{in} 都是一次端口周期，CPU 在 \code{.wait} 上空转的拍数就是查询开销：

\begin{lstlisting}
; bl = 待发送字节
mov dx, 0x3FD       ; LSR
.wait:
in  al, dx          ; 读状态寄存器 = 一次状态采样
test al, 0x20       ; THRE: 发送保持寄存器空?
jz  .wait
mov dx, 0x3F8       ; THR
mov al, bl
out dx, al
\end{lstlisting}

LSR 的关键位：

\begin{longtable}{L{0.10\linewidth}L{0.16\linewidth}L{0.66\linewidth}}
\toprule
位 & 名称 & 含义 \\
\midrule
0 & DR & 接收数据就绪：RBR 里有未读字节（读 RBR 后由硬件清零） \\\tablerowrule
1 & OE & 溢出错误：新字节到达时 RBR 还没被读走，旧字节被覆盖 \\\tablerowrule
5 & THRE & 发送保持寄存器空：可以接受下一个待发送字节 \\\tablerowrule
6 & TEMT & 发送器完全空：移位寄存器也空了，最后一位已真正上线 \\
\bottomrule
\end{longtable}

16550 相对早期 8250 的主要改进是 16 字节收发 FIFO，而 FCR 的触发水位（1/4/8/14 字节可选）体现的是同一个权衡：水位设 14，则收满 14 字节才发一次中断，中断频率降为逐字节中断的十几分之一，但软件看到数据的延迟最长可达 14 个字节时间。它就是 1.3 节“查询周期 $P$”在硬件里的翻版，只是这个环由硬件闭合，软件只在配置时选一次参数。

\section{五、微序列器：控制字怎样按条件推进}

把复杂指令写成多行微码，只解决了“每一行输出哪些控制线”，还没有回答“下一行从哪里来”。微序列器（microsequencer）保存微程序计数器 \(\mu PC\)，用当前微指令的下一地址字段、条件码、等待输入和主指令入口表共同选择下一 \(\mu PC\)。控制存储器读出的控制字一部分直接驱动数据通路，另一部分回到微序列器，决定顺序前进、条件分支、调用公共子程序、返回或结束本条指令。

\begin{pdfdiagram}[x=1cm,y=1cm]
  \node[hw memory, minimum width=2.4cm] (ir) at (1.3,4.0) {指令寄存器\\opcode 与模式};
  \node[hw control, minimum width=2.5cm] (entry) at (4.2,4.0) {入口映射\\首条微地址};
  \node[hw state, minimum width=2.0cm] (upc) at (7.0,4.0) {\(\mu PC\)};
  \node[hw memory, minimum width=2.6cm] (store) at (10.1,4.0) {控制存储器\\控制字与下一地址};
  \node[hw compute, minimum width=2.6cm] (datapath) at (10.1,1.3) {数据通路动作\\读、运算、写候选};
  \node[hw control, minimum width=2.8cm] (next) at (6.8,1.3) {下一地址选择\\条件/等待/返回};
  \node[hw state, minimum width=2.5cm] (inputs) at (3.5,1.3) {状态输入\\标志、ready、异常};
  \draw[hw bus] (ir) -- (entry);
  \draw[hw bus] (entry) -- (upc);
  \draw[hw bus] (upc) -- (store);
  \draw[hw bus] (store) -- (datapath);
  \draw[hw bus] (store) -- (next);
  \draw[hw bus] (inputs) -- (next);
  \draw[hw return] (next) -- (upc);
\end{pdfdiagram}
\diagramnote{微序列器的闭环。主指令通过入口映射给出首条微地址；控制存储器同时输出数据通路控制和下一地址信息；标志、ready 与异常参与下一地址选择。只有结束微指令到达提交条件后，主指令才完成，\(\mu PC\) 的中间推进不是架构完成。}

\topic{等待状态必须保持控制幂等，不能重复外部副作用}

若某条微指令启动存储器读，控制器可能要在 \code{ready=0} 时停留数拍。最安全的组织是把“发出一次请求”和“等待完成”分成不同微状态：发出状态只在握手成功的那个周期转移所有权，等待状态只观察完成条件，不再次拉起会被从设备当成新请求的有效信号。若协议采用 valid/ready，valid 可以保持，但请求内容和事务身份必须稳定；接收一次 ready 后立即转入下一状态。

同理，写寄存器、递增 PC 和写 MMIO 不能放在会反复执行的等待控制字中。内部暂存寄存器可以在每拍覆盖同一候选值，架构寄存器和外部设备却必须只更新一次。微码审查不仅要看每行控制位，还要标出每个动作是“可重复计算”“握手一次”“架构提交”还是“不可重放副作用”。

\begin{longtable}{L{0.18\linewidth}L{0.28\linewidth}L{0.27\linewidth}L{0.17\linewidth}}
\toprule
微状态类型 & 允许重复的动作 & 离开条件 & 错误表现 \\
\midrule
入口/译码 & 组合形成候选地址 & 指令和模式合法 & 进入错误微程序 \\\tablerowrule
请求发出 & 保持地址、数据和身份稳定 & valid 与 ready 同拍成立 & 重复事务或请求丢失 \\\tablerowrule
等待完成 & 读取状态，保持暂存值 & completion 与身份匹配 & 把旧完成认作当前完成 \\\tablerowrule
内部运算 & 覆盖私有临时寄存器 & 运算状态稳定或计数结束 & 中间值提前可见 \\\tablerowrule
架构提交 & 一次写入寄存器、PC 或标志 & 无更高优先级异常 & 重复副作用或不精确异常 \\
\bottomrule
\end{longtable}

\topic{公共微子程序节省控制存储，但要保存有限返回状态}

地址形成、规格化和权限检查可能被多条指令复用。微序列器可以提供一层或少量深度的返回栈：调用微指令保存 \(\mu PC+1\)，跳到公共入口；返回微指令恢复保存地址。这个栈不是软件可见调用栈，容量通常固定，溢出应在设计验证中被证明不可达。若公共序列可能等待、异常或再调用，返回地址、原指令身份和中间状态都必须保持到它结束。

\section{六、可重启微操作与微码补丁：内部复杂度怎样不破坏架构边界}

一条架构指令可以展开成多条内部微操作，但软件通常只允许看到“未执行”或“已完成”两个状态。处理器因此把地址、部分结果和循环计数保存在不可见临时状态中，直到全部检查通过才一次更新架构寄存器、标志和 PC。若中途发生可处理缺页或外部中断，控制器可以丢弃临时状态并从原指令重启，或在体系结构明确允许的检查点保存内部进度；两种方案都必须保证已对外产生的副作用不会被重复。

\begin{pdfdiagram}[x=1cm,y=1cm]
  \node[hw state, minimum width=2.5cm] (start) at (1.3,4.0) {架构旧状态\\PC、寄存器、标志};
  \node[hw control, minimum width=2.5cm] (expand) at (4.2,4.0) {译码/assist 入口\\建立微操作身份};
  \node[hw memory, minimum width=2.5cm] (temp) at (7.1,4.0) {临时状态\\地址、部分结果、计数};
  \node[hw control, minimum width=2.5cm] (check) at (10.0,4.0) {最终检查\\异常与完成条件};
  \node[hw state, minimum width=2.5cm] (commit) at (10.0,1.3) {架构提交\\结果与后继 PC};
  \node[hw control, minimum width=2.5cm] (recover) at (7.1,1.3) {取消/恢复\\清理临时所有权};
  \node[hw state, minimum width=2.5cm] (retry) at (4.2,1.3) {保存原因\\原 PC 重试};
  \draw[hw bus] (start) -- (expand);
  \draw[hw bus] (expand) -- (temp);
  \draw[hw bus] (temp) -- (check);
  \draw[hw bus] (check) -- (commit);
  \draw[hw return] (check) -- (recover);
  \draw[hw return] (recover) -- (retry);
  \draw[hw return] (retry) -- (start);
\end{pdfdiagram}
\diagramnote{复杂指令的内部状态与架构状态分离。上排可以执行许多微操作并暂存部分结果；只有最终检查通过才沿右侧提交。异常路径先撤销临时所有权，再以原 PC 和确定原因进入重试，不能把半完成结果留给软件。}

\topic{内存和 MMIO 副作用决定一条微序列能否整体重启}

纯寄存器计算只要结果尚未提交，通常可以整体丢弃并重算。读普通内存也可在权限和一致性规则允许时重发，因为重复读取不会写入外部状态。写内存则需要在提交点之前保存在 store buffer 或私有写队列，待整条指令无异常后再获得可见性；若微序列直接写到 Cache 或总线后才发现异常，整体重启可能重复写入。

MMIO 更严格。读设备状态可能具有 read-to-clear 语义，写门铃可能立即启动外部动作，这些事务不能假设可重放。实现要么在所有可能故障检查完成后才发出 MMIO，要么把指令定义成发生该副作用后不再报告可重试故障。可重启性不是“保存一个 \(\mu PC\)”就完成，而取决于微序列之前是否已经越过不可撤销边界。

\topic{微码 assist 和补丁仍要经过版本、状态与回滚验证}

现代核心会把少见复杂指令、特殊输入或硬件勘误处理转入微码 assist。assist 注入的内部操作可以比常规译码路径更长，因此会改变延迟和吞吐，却不能改变 ISA 规定的结果、异常优先级和提交语义。性能计数器若显示 assist 增多，说明工作负载频繁进入慢路径；这是一条定位证据，不等于数据已经错误。

可更新微码通常在复位或受控更新流程中装入受验证的补丁映像。平台要核对处理器签名、补丁版本和认证结果，再使新入口或替换序列生效；多核系统还要保证所有目标核心进入一致版本。更新失败时应保留原有可启动路径，不能让一部分核心运行旧序列、另一部分核心运行新序列。补丁改变的是控制实现，架构状态格式和软件可见契约必须保持兼容。

\begin{longtable}{L{0.18\linewidth}L{0.27\linewidth}L{0.28\linewidth}L{0.17\linewidth}}
\toprule
边界 & 内部状态处理 & 架构可见结果 & 恢复要求 \\
\midrule
普通完成 & 临时结果通过全部检查 & 一次更新目的状态与后继 PC & 释放微序列资源 \\\tablerowrule
可重试故障 & 丢弃或保存规定检查点 & 目的状态不变，保存原指令 PC & 修复后执行语义等价 \\\tablerowrule
不可重放副作用 & 在发出前完成全部可故障检查 & 外部动作只发生一次 & 禁止从副作用之前盲目重启 \\\tablerowrule
微码更新 & 验证映像、目标和版本后切换 & ISA 契约保持不变 & 失败回到已验证版本 \\
\bottomrule
\end{longtable}

\section*{章末练习}

\begin{chapterexercise}{寄存器堆}
说明读地址、写地址、写数据、写使能如何配合完成 \code{R3=R1+R2}
\exercisecheck{完成要求}{写回发生在时钟沿}
\end{chapterexercise}

\begin{chapterexercise}{4 寄存器实验}
用 74HC173/574、74HC138、mux 设计 4 个 4-bit 寄存器的读写方案
\exercisecheck{完成要求}{同一周期只写一个目标，A/B 两个读端口可独立选择}
\end{chapterexercise}

\begin{chapterexercise}{总线冲突}
解释为什么 ALU 输出和存储器输出不能随意直接并联
\exercisecheck{完成要求}{能说明 mux 或三态使能的必要性}
\end{chapterexercise}

\begin{chapterexercise}{控制器}
比较硬连线控制、状态机控制和微码控制
\exercisecheck{完成要求}{能说明适用场景和风险}
\end{chapterexercise}

\begin{chapterexercise}{控制 ROM}
给出一个控制 ROM 地址格式和控制字字段，说明 ADD 与条件分支各看哪些输入
\exercisecheck{完成要求}{每个控制字位都能追到实际控制线}
\end{chapterexercise}

\begin{chapterexercise}{控制 trace}
为三周期 \code{R3=R1+R2} 列出 S0/S1/S2 的打开路径和提交状态
\exercisecheck{完成要求}{能定位每个周期的写使能}
\end{chapterexercise}

\begin{chapterexercise}{最小 CPU 接口}
列出 PC、IR、寄存器堆、ALU、存储器接口和控制器之间的关键输入输出
\exercisecheck{完成要求}{能说明哪些是组合路径，哪些在时钟沿提交}
\end{chapterexercise}

\begin{chapterexercise}{错误定位}
给出 \code{alu_op}、\code{wb_sel}、\code{reg_write} 错误各自导致的现象
\exercisecheck{完成要求}{能按 trace 排查，而不是盲改 ALU}
\end{chapterexercise}

\begin{chapterexercise}{微程序读法}
说明由 opcode、step 和标志拼成的地址如何读出控制字并驱动控制线
\exercisecheck{完成要求}{能把查表动作和实际电线对应起来}
\end{chapterexercise}

\section*{参考解答与要点}

\begin{chapteranswer}{寄存器堆}
\code{read_a=R1, read_b=R2, a_sel=reg_a, b_sel=reg_b, alu_op=ADD, wb_sel=ALU, write_addr=R3, reg_write=1}。若采用多周期，还需 A/B 和 ALUOut 写使能。
\answercheck{必须答到的要点}{写回必须等 ALU 输出稳定并在时钟沿提交。}
\end{chapteranswer}

\begin{chapteranswer}{4 寄存器实验}
四个寄存器保存 R0 到 R3；写地址经译码产生四路候选写使能，再与 \code{reg_write} 相与；读端口 A/B 分别用 mux 按读地址选择一路输出；写回数据并接到所有寄存器 D 端但只有被选寄存器写入。
\answercheck{必须答到的要点}{读是组合选择，写是时钟沿提交。}
\end{chapteranswer}

\begin{chapteranswer}{总线冲突}
两个输出若同时驱动同一线且电平不同，会形成冲突和大电流。应使用 mux 明确选择一路，或用三态总线并保证同一时刻只有一个输出使能有效。
\answercheck{必须答到的要点}{共享连线必须有唯一驱动者。}
\end{chapteranswer}

\begin{chapteranswer}{控制器}
硬连线控制适合简单高频路径；状态机适合多周期阶段；微码适合复杂指令分解；混合方案让常见路径快、复杂路径可管理。
\answercheck{必须答到的要点}{所有方案最终都要输出具体控制线。}
\end{chapteranswer}

\begin{chapteranswer}{控制 ROM}
地址可由 \code{opcode}、微步骤 \code{step}、标志位和复位/中断请求组成；数据输出可包含 \code{alu_op}、\code{a_sel}、\code{b_sel}、\code{wb_sel}、\code{reg_write}、\code{flags_write}、\code{pc_write}、\code{ir_write}、\code{mem_read/write}、\code{next_sel}。ADD 看 opcode 与 step，条件分支还要看 Zero 等标志。
\answercheck{必须答到的要点}{ROM 每个输出位都必须连接到一个真实控制对象。}
\end{chapteranswer}

\begin{chapteranswer}{控制 trace}
S0 打开 R1/R2 到 A/B 的路径并提交 A/B；S1 打开 A/B 到 ALU 的路径并提交 ALUOut/Flags；S2 选择 ALUOut 写回 R3 并打开 \code{reg_write}。
\answercheck{必须答到的要点}{每周期要写清打开路径、提交状态和关闭信号。}
\end{chapteranswer}

\begin{chapteranswer}{最小 CPU 接口}
PC 输出取指地址并由 \code{pc_write/pc_sel} 更新；IR 锁存指令并输出 opcode/寄存器字段；寄存器堆输出操作数并接收写回；ALU 产生结果和标志；存储器接口处理地址、读写数据和等待；控制器输入 opcode、状态、标志和等待信号并输出控制字。
\answercheck{必须答到的要点}{PC、IR、寄存器堆和标志是状态元件，必须有写使能。}
\end{chapteranswer}

\begin{chapteranswer}{错误定位}
\code{alu_op} 错会算错函数；\code{wb_sel} 错会写回错误来源；\code{reg_write} 错会漏写或误写寄存器。应按控制 trace 检查状态、选择线和写使能。
\answercheck{必须答到的要点}{先区分计算错误、选择错误和提交错误。}
\end{chapteranswer}

\begin{chapteranswer}{微程序读法}
把 opcode、微步骤 step 和要检查的标志拼成 ROM 地址；该地址读出的数据行就是本周期控制字，各字段直接接 \code{alu_op}、\code{a_sel}、\code{wb_sel}、\code{reg_write}、\code{next_sel} 等控制线。例如条件分支在 EXEC 步且 Z=1 时，控制字让 \code{next_sel} 选择分支目标而非顺序状态。
\answercheck{必须答到的要点}{微指令不是被执行的软件语句，而是被查出的电平表。}
\end{chapteranswer}
